\documentclass[12pt]{article}  
\usepackage{amsmath}
\usepackage{amssymb}

\usepackage{booktabs}
\usepackage{array}

\usepackage{hyperref}
\providecommand{\href}[1]{\url{#1}}
\usepackage{iftex}
\usepackage[utf8]{inputenc}  
\usepackage{graphicx}
\setlength{\parindent}{0pt}
\setlength{\parskip}{6pt plus 2pt minus 1pt}
\usepackage[normalem]{ulem}
\usepackage[
]{geometry}

\usepackage{amsthm}
\usepackage{float}

\theoremstyle{plain}
\newtheorem{theorem}{Theorem}[section]
\newtheorem{lemma}{Lemma}
\newtheorem{corollary}{Corollary}
\newtheorem{proposition}{Proposition}
\newtheorem{conjecture}{Conjecture}

\theoremstyle{remark}
\newtheorem{remark}{Remark}
\newtheorem{note}{Note}
\newtheorem{claim}{Claim}

\theoremstyle{definition}
\newtheorem{definition}{Definition}[section]  
\newtheorem{condition}{Condition}
\newtheorem{problem}{Problem}[section]
\newtheorem{example}{Example}[section]

\renewcommand\qedsymbol{$\blacksquare$}  

\usepackage{mathtools}
\usepackage{listings}

\usepackage{framed}

\ProvidesPackage{iftex}

\ifxetex

\usepackage{unicode-math}
% \setmainfont{STIX Two Text}
% \setmathfont{STIX Two Math}
% \setmathfont[range={\mathcal,\mathbfcal},StylisticSet=1]{STIX Two Math}

% \setmonofont{JetBrains Mono}

\fi

\providecommand{\tightlist}{\setlength{\itemsep}{0pt}\setlength{\parskip}{0pt}} 

\usepackage{csquotes}
\title{Linear Algebra}
\author{https://github.com/icprplshelp/}
\date{\today}

\begin{document}


\maketitle

{
\setcounter{tocdepth}{3}
\tableofcontents
}

\section{Linear transformations}

A function \(T:\mathbb{R}^{n} \rightarrow \mathbb{R}^{m}\) is a linear
transformation if
\(\forall\mathbf{v},\ \mathbf{u} \in \mathbb{R}^{n}\) and
\(\forall r\mathbb{\in R}\) and has the following properties:

\begin{enumerate}
\def\labelenumi{(\roman{enumi})}
\item
  \(T\left( \mathbf{u} + \mathbf{v} \right) = T\left( \mathbf{u} \right) + T(\mathbf{v})\)
\item
  \(T\left( r\mathbf{v} \right) = rT(\mathbf{v})\)
\end{enumerate}

Example: \(f(x) = kx\) is a linear function
\(\mathbb{R}^{1} \rightarrow \mathbb{R}^{1}\). The following properties
hold:

\begin{align*}
f\left( x_{1} + x_{2} \right) &= k\left( x_{1} + x_{2} \right) = f\left( x_{1} \right) + f\left( x_{2} \right) \\
f\left( rx_{1} \right) &= k\left( rx_{1} \right) = rkx_{1} = rf(x_{1})
\end{align*}


Second example: \(f(x) = x^{2}\) is not a linear function
\(\mathbb{R}^{1} \rightarrow \mathbb{R}^{1}\) a.k.a.

\[f(2x) \neq 2f(x)
\]

RHS: \(2f(x) \Rightarrow \text{LHS} \neq \text{RHS}\)

\begin{remark}
To show that a function is linear, you can combine (i)
and (ii) together.

\[{(i) + \left( ii \right) = f\left( \mathbf{u} + r\mathbf{v} \right) }{= f\left( \mathbf{u} \right) + rf\left( \mathbf{v} \right)\ \forall\mathbf{u},\ \mathbf{v} }{\in \mathbb{R}^{n},\ \forall r\mathbb{\in R}} \]
\end{remark}

\textbf{Capital letters may be used to denote a linear function. Don't
confuse them and matrices!}


\subsection{What are functions?}

\textbf{A function \emph{compacts} complex statements. For example:}

Let \(h(x) = 4x^{2} + 6x\)

Then we can write \(4 + h(x)\) as \(4 + 4x^{2} + 6x\), and anytime we
see \(4x^{2} + 6x\) appear in a statement, we can convert it back to
\(h(x)\). Note that \(h(x)\) doesn't have to be represented in function
notation; if we let \(y = h(x)\), then we could write \(4 + h(x)\) as
\(4 + y\). Knowing this is essential to prevent misunderstandings. That
being said, if we define \(g(x) = 1\) then whenever \(g(x)\) appears
anywhere in a statement we can replace it with \(1\) (isn't that just a
constant?).


\subsection{What are sets?}

Sets are an unordered collection of anything. For example,
\(\{ 1,\ 2,\ 3,\ 4\}\) is a set. We can also have a set of real numbers,
or a set of even numbers, which would be denoted as
\(\left\{ x\mathbb{\in Z}\mid x\% 2 = 0 \right\}\). That's set builder
notation, and here's how it works: \{our set \textbar{} if it meets what
we write here\}. Knowing this will help mitigate confusion further on.


\subsection{Matrices as a form of a linear
function}

If these properties hold, then \(f(\mathbf{x})\) is linear. \(A\) is
a matrix. Remember that \(A\mathbf{x}\) returns a vector.

\begin{align*}
f\left( \mathbf{x} \right) &= A\mathbf{x} = \mathbf{y} \Rightarrow f:\mathbb{R}^{n} \rightarrow \mathbb{R}^{m} \\
A\left( \mathbf{x} + \mathbf{y} \right) &= A\mathbf{x} + A\mathbf{y} \\
A\left( r\mathbf{x} \right) &= rA\mathbf{x},\ r\mathbb{\in R} \\
&\Rightarrow f\left( \mathbf{x} \right)\text{\ is\ linear}
\end{align*}



\subsection{Domain, Codomain, Image, and
Kernel}

A linear transformation's domain, codomain, image, and kernel are sets,
and thus can be described in the set builder notation.

\textbf{Definition:} If \(T:\mathbb{R}^{n} \rightarrow \mathbb{R}^{m}\)
is a linear transformation then\ldots{}

\begin{enumerate}
\def\labelenumi{\arabic{enumi}.}
\item
  \(\mathbb{R}^{n}\) is the domain of \(T\). (The set of inputs)
\item
  \(\mathbb{R}^{m}\) is the codomain of \(T\). (The set of possible
  outputs)
\item
  Let \(w \subseteq \mathbb{R}^{n}\), then the image of \(w\) under
  \(T\) is
  \(T(w) = \left\{ T\left( \mathbf{w} \right) \middle| \mathbf{w} \in W \right\}\)
\end{enumerate}


\begin{enumerate}
\def\labelenumi{\alph{enumi}.}
\item
  \(T(w)\) is a collection of all possible outputs corresponding to the
  inputs of \(w\). \(T(w)\) is the image.
\end{enumerate}

\begin{enumerate}
\def\labelenumi{\arabic{enumi}.}
\setcounter{enumi}{3}
\item
  The range of \(T\) is
  \(T\left( \mathbb{R}^{n} \right) = \left\{ T\left( \mathbf{v} \right) \middle| \mathbf{v} \in \mathbb{R}^{n} \right\}\)
  (i.e., Range of \(T\) is the image of the domain.)
\item
  Let \(w' \subseteq \mathbb{R}^{m}\). The inverse image of \(w'\)
  under \(T\) is
  \(T^{- 1}\left( w' \right) = \left\{ \mathbf{v} \in \mathbb{R}^{n} \middle| T\left( \mathbf{v} \right) \in w' \right\}\)
  (is this the domain?)
\item
  Let \(\mathbf{0}\) be the zero vector in \(\mathbb{R}^{m}\).
\end{enumerate}

\(T^{- 1}\left( \mathbf{0} \right) = \left\{ \mathbf{v} \in \mathbb{R}^{n} \middle| T\left( \mathbf{v} \right) = \mathbf{0} \in \mathbb{R}^{m} \right\}\)
is called the kernel of \(T\). \textbf{(The \emph{kernel} of a linear
transformation \emph{T} is the set of inputs that \emph{T} sends to the
zero vector.)}


\subsection{An example of a linear
transformation}

\textbf{Example:} Let \(T:\mathbb{R}^{3} \rightarrow \mathbb{R}^{2}\).
\(T\left( \left\lbrack x_{1},\ x_{2},\ x_{3} \right\rbrack \right) = \lbrack x_{1} - x_{2},\ x_{2} + x_{3}\rbrack\).

This is one way for us to define a function and is a definition of a
function.

\textbf{Here's a series of questions:}

\begin{enumerate}
\def\labelenumi{\arabic{enumi}.}
\item
  \textbf{DETERMINING WHETHER \emph{T} IS LINEAR}
\end{enumerate}

Is \(T\) a linear map/function/transformation? \textbf{Answer: Yes.
Why?}

\(T\left( \mathbf{u} + r\mathbf{v} \right) = T\left( \mathbf{u} \right) + rT\left( \mathbf{v} \right)\ \forall\mathbf{u},\mathbf{v},\ \ r\mathbb{\in R}\).

If we verify if LHS of the above equation is always equal to RHS,
\textbf{then} \(T\) is a linear function. Our objective is to start with
\(T(\mathbf{u} + r\mathbf{v})\) and perform some operations to
get us to

\[T\left( \mathbf{u} \right) + rT(\mathbf{v})
\]

\[r\mathbb{\in R}\]

\emph{Rationale:} \(\mathbf{u}\) \emph{and} \(\mathbf{v}\)
\emph{are inputs to} \(T\) \emph{and must be} \(\mathbb{R}^{3}\)\emph{.}

LHS:

\begin{align*}
T\left( \mathbf{u} + r\mathbf{v} \right) &= T\left( \left\lbrack u_{1},\ u_{2},\ u_{3} \right\rbrack + r\left\lbrack v_{1},\ v_{2},\ v_{3} \right\rbrack \right) \\
&= T\left( \left\lbrack u_{1} + rv_{1},\ u_{2} + rv_{2},\ u_{3} + rv_{3} \right\rbrack \right)
\end{align*}


Running the function, we get this. We can do some algebraic manipulation
to get what we have as close to the RHS,

\begin{align*}
&= \left\lbrack u_{1} + rv_{1} - u_{3} - rv_{3},\ u_{2} + rv_{2} + u_{3} + rv_{3} \right\rbrack \\
&= \left\lbrack u_{1} - u_{3} + r\left( v_{1} - v_{3} \right),\ u_{2} + u_{3} + r\left( v_{2} + v_{3} \right) \right\rbrack \\
&= \left\lbrack u_{1} - u_{3},\ u_{2} + u_{3} \right\rbrack + \lbrack r\left( v_{1} - v_{3} \right),\ r\left( v_{2} + v_{3} \right)\rbrack \\
&= \left\lbrack u_{1} - u_{3},\ u_{2} + u_{3} \right\rbrack + r\lbrack v_{1} - v_{3},\ v_{2} + v_{3}\rbrack
\end{align*}


Recall:

\begin{align*}
T\left( x_{1},x_{2},\ x_{3} \right) &= \left( x_{1} - x_{3},\ x_{2} + x_{3} \right) \\
&= T\left( \left\lbrack u_{1},\ u_{2},\ u_{3} \right\rbrack \right) + rT\left( \left\lbrack v_{1},\ v_{2},\ v_{3} \right\rbrack \right)
\end{align*}


\[T\left( \mathbf{u} \right) + rT\left( \mathbf{v} \right) = RHS\]

Therefore \(LHS = RHS \Rightarrow T\) is a linear transformation.

\begin{enumerate}
\def\labelenumi{\arabic{enumi}.}
\setcounter{enumi}{1}
\item
  \textbf{FINDING THE IMAGE OF A FUNCTION WITH AN INPUT DOMAIN}
\end{enumerate}

Let
\(H = \left\{ \lbrack x,\ x,\ x\rbrack \in \mathbb{R}^{3} \right\} \subseteq \mathbb{R}^{3}\).
Find the \textbf{image} of \(H\) under \(T\).

\textbf{Recall:} We let \(T:\mathbb{R}^{3} \rightarrow \mathbb{R}^{2}\).
\(T\left( \left\lbrack x_{1},\ x_{2},\ x_{3} \right\rbrack \right) = \lbrack x_{1} - x_{2},\ x_{2} + x_{3}\rbrack\).

\textbf{Objective:} Want to find the image of \emph{H} under \emph{T} in
the most explicit way as possible.

\[T(H) = \left\{ \mathbf{y} \in \mathbb{R}^{2} \right|\mathbf{y} = T\left( \mathbf{u} \right),\ \mathbf{u} \in H\}\]

The above notation says \(\mathbf{y}\) is in \(\mathbb{R}^{2}\)
(because \(T(H)\) is in \(\mathbb{R}^{2}\)), but only if
\(\mathbf{y}\) is an output of \(T(\mathbf{u})\) where
\(\mathbf{u}\) is in \(H\), a set of \(\mathbb{R}^{3}\) vectors
where all its entries are identical. We \emph{could} stop at this step,
but why not expand this further?

\begin{align*}
&= \left\{ \mathbf{y} \in \mathbb{R}^{2} \middle| \mathbf{y} = T\left( \lbrack u,\ u,\ u\rbrack \right),\ \mathbf{u} = \lbrack u,\ u,\ u\rbrack \in H \right\} \\
&= \{\mathbf{y} \in \mathbb{R}^{2}|\mathbf{y} = \lbrack 0,\ u + u\rbrack = \lbrack 0,\ 2u\rbrack,\ u\mathbb{\in R\}}
\end{align*}


Conclude that
\(\left\{ \mathbf{y} \in \mathbb{R}^{2} \middle| \mathbf{y} = \lbrack 0,\ 2u\rbrack,\ u\mathbb{\in R} \right\}\)

This is the image of \(H\) under \(T\). This is an explicit
representation of an image without \(T(\mathbf{u})\) appearing
anywhere within the set builder.

\begin{enumerate}
\def\labelenumi{\arabic{enumi}.}
\setcounter{enumi}{2}
\item
  \textbf{DEDUCING INVERSE IMAGES WITH AN IMAGE}
\end{enumerate}

\textbf{Question:} Let
\(u = \left\{ \lbrack 1,\ 2\rbrack,\ \lbrack - 1,\ 3\rbrack \right\} \subseteq \mathbb{R}^{2}\)
(which is a subset of the codomain of \(T\). Find the inverse image of
\(u\), under \(T\).

We are given a \textbf{set} \(u\) with two elements -- it's like two
dots in a set, so what we actually have to do is compute a function's
inverse with these inputs -- we were given two inputs, which means
double the work!

\textbf{Recall:} We let \(T:\mathbb{R}^{3} \rightarrow \mathbb{R}^{2}\).
\(T\left( \left\lbrack x_{1},\ x_{2},\ x_{3} \right\rbrack \right) = \lbrack x_{1} - x_{2},\ x_{2} + x_{3}\rbrack\).

The set builder below denotes our inverse image. We just have to compact
that down, so we end up with an actual set.

\begin{align*}
T^{- 1}(u) &= \left\{ \mathbf{x} = \left\lbrack x_{1},\ x_{2},\ x_{3} \right\rbrack \in \mathbb{R}^{3} \right|\ T\left( \mathbf{x} \right) \in u\} \\
\Rightarrow T\left( \mathbf{x} \right) &= \lbrack 1,\ 2\rbrack\ or\ \lbrack - 1,\ 3\rbrack
\end{align*}


\textbf{Case 1:}
\(T\left( \mathbf{x} \right) = \lbrack 1,\ 2\rbrack\). You'll need
to find \(\lbrack x_{1},\ x_{2},\ x_{3}\rbrack\) such that
\(T\left( \left\lbrack x_{1},x_{2},x_{3} \right\rbrack \right) = \lbrack 1,\ 2\rbrack\)

\begin{align*}
\left\lbrack x_{1} - x_{3},\ x - 2 + x_{3} \right\rbrack &= \lbrack 1,\ 2\rbrack \Rightarrow \begin{matrix}x_{1} - x_{3} = 1 \\x_{2} + x_{3} = 2 \\\end{matrix} \\
&\Rightarrow \left\lbrack \begin{matrix}1 & 0 & - 1 \\0 & 1 & 1 \\\end{matrix}\  \middle| \ \begin{matrix}1 \\2 \\\end{matrix} \right\rbrack
\end{align*}


\emph{So, you might be asking, where did the augmented matrix come
from?} That was because a system of linear equations appeared in your
work, so you must solve it. It is a 3-unknown 2-equation system, so you
will end up with a free variable.

3-variable, 2-equation systems typically have infinite solutions, and
can be solved like any other linear system. If you're having trouble
making sense out of this, you can always add a useless equation
\(0x_{1} + 0x_{2} + 0x_{3} = 0\) giving you

\[\left\lbrack \begin{matrix}
1 & 0 & - 1 \\
0 & 1 & 1 \\
0 & 0 & 0 \\
\end{matrix}\, \middle| \,\begin{matrix}
1 \\
2 \\
0 \\
\end{matrix} \right\rbrack\]

Which will have a free variable, resulting in infinitely many solutions.

This will give you: \(\left\{ \begin{matrix}
x_{1} - x_{3} = 1 \\
x_{2} + x_{3} = 2 \\
\end{matrix} \right.\  \Rightarrow \begin{matrix}
x_{3} = s \\
x_{1} = 1 + s \\
x_{2} = 2 - s \\
\end{matrix}\)

We can thus conclude that:

\[\mathbf{x} = \left\lbrack x_{1},x_{2},\ x_{3} \right\rbrack = \lbrack 1 + s,\ 2 - s,\ s\rbrack,\ s\mathbb{\in R}\]

We deduced a part of the inverse image of \(u\) under \(T\).

\[T\left( \left\lbrack \frac{1 + s}{x_{1}},\frac{2 - s}{x_{2}},\frac{s}{x_{3}} \right\rbrack \right) = \left\lbrack \frac{1 + s}{x_{1}} - \frac{s}{x_{3}},\frac{2 - s}{x_{2}} + \frac{s}{x_{3}} \right\rbrack = \lbrack 1,\ 2\rbrack\]

\textbf{Case 2:}
\(T\left( \mathbf{x} \right) = \left\lbrack x_{1} - x_{3},\ x_{2} + x_{3} \right\rbrack = \lbrack - 1,\ 3\rbrack\)

\[\Rightarrow \left\{ \begin{matrix}
x_{1} - x_{3} = - 1 \\
x_{2} + x_{3} = 3 \\
\end{matrix} \right.\  \Rightarrow \mathbf{x} = \lbrack - 1 + t,\ 3 - t,\ t\rbrack,\ t\mathbb{\in R}\]

That's case 2, complete. Now we can combine case 1 and 2 together and
that is our inverse image, which is:

\[\lbrack 1 + s,\ 2 - s,\ s\rbrack\ \text{and}\ \lbrack - 1 + t,\ 3 - t,\ t\rbrack\]

Or we can represent it as a set:
\(\left\{ \lbrack 1 + s,\ 2 - s,\ s\rbrack,\ \lbrack - 1 + t,\ 3 - t,\ t\rbrack \right\}\)

\textbf{Because} \(\mathbf{u}\) \textbf{is an ``image'', we can
represent the inverse image like the following:}

\[T\left( \left\{ \lbrack 1 + s,\ 2 - s,\ s\rbrack,\ \lbrack - 1 + t,\ 3 - t,\ t\rbrack\mid t,\ s\mathbb{\in R} \right\} \right) = u\]

\[T\left( \left\{ \begin{bmatrix}
1 + s \\
2 - s \\
s \\
\end{bmatrix},\ \begin{bmatrix}
 - 1 + t \\
3 - t \\
t \\
\end{bmatrix}\mid t,\ s\mathbb{\in R} \right\} \right) = u\]

\begin{enumerate}
\def\labelenumi{\arabic{enumi}.}
\setcounter{enumi}{3}
\item
  \textbf{FINDING THE KERNEL OF A FUNCTION}
\end{enumerate}

Find the kernel of \emph{T}.

\[\ker(T) = \{\mathbf{x} \in \mathbb{R}^{3}|T\left( \mathbf{x} \right) = 0\}\]

Figure out when input \(\mathbf{x}\) can give you 0 as its output.

\textbf{Recall:}
\(T\left( \mathbf{x} \right) = T\left( \left\lbrack x_{1},x_{2},\ x_{3} \right\rbrack \right) = \left\lbrack x_{1} - x_{3},\ x_{2} + x_{3} \right\rbrack = \lbrack 0,\ 0\rbrack\)

This can be used to build a system of equations.

\begin{align*}
&\left\lbrack \begin{matrix}1 & 0 & - 1 \\0 & 1 & 1 \\\end{matrix}\, \middle| \,\begin{matrix}0 \\0 \\\end{matrix} \right\rbrack \Rightarrow \lbrack s,\  - s,\ s\rbrack,\ s\mathbb{\in R} \\
T\left( \lbrack s,\  - s,\ s\rbrack \right) &= \lbrack s - s,\  - s + s\rbrack = \lbrack 0,\ 0\rbrack \Rightarrow \mathbf{0} \in \mathbb{R}^{2}
\end{align*}


\[\ker{(T)} = \left\{ \lbrack s,\  - s,\ s\rbrack \in \mathbb{R}^{3} \middle| s\mathbb{\in R} \right\}\]


\subsection{Basis vectors}

\textbf{Proposition:} Let
\(T:\mathbb{R}^{n} \rightarrow \mathbb{R}^{m}\) be a linear
transformation. Then

\begin{align*}
&T\left( r_{1}{\mathbf{v}}_{1} + r_{2}{\mathbf{v}}_{2} + \ldots + r_{k}{\mathbf{v}}_{k} \right) \\
&= rT\left( {\mathbf{v}}_{1} \right) + r_{2}T\left( {\mathbf{v}}_{2} \right) + \cdots + r_{k}T\left( {\mathbf{v}}_{k} \right)
\end{align*}


This is a linear combination between
\(\mathbf{v_{1}},\ \cdots,\ {\mathbf{v}}_{k}\)

\textbf{Remark:} What if \(\mathbf{v_{1}}\cdots\mathbf{v_{n}}\)
is a set of \textbf{basis} vectors of \(\mathbb{R}^{n}\), i.e.,
\(\text{span}\left( \mathbf{v_{1}}\cdots\mathbf{v_{n}} \right) = \mathbb{R}^{n}\)
and \({\mathbf{v}}_{1}\cdots{\mathbf{v}}_{n}\) are linearly
independent

Then

\[\Leftrightarrow \begin{bmatrix}
| & | & | \\
{\mathbf{v}}_{1} & \cdots & {\mathbf{v}}_{n} \\
| & | & | \\
\end{bmatrix}\sim\begin{bmatrix}
1 & \cdots & 0 \\
 \vdots & \ddots & \vdots \\
0 & \cdots & 1 \\
\end{bmatrix}\]

\textbf{Proposition:} For any input
\(\mathbf{\alpha} \in \mathbb{R}^{n}\text{.\ }\mathbf{\alpha} = r_{1}\mathbf{v_{1}} + r\mathbf{v_{2}} + \cdots + r\mathbf{v_{n}}\)
(given \(\mathbf{v_{1}}\cdots\mathbf{v_{n}}\) are a set of basis
vectors)

\[T\left( \mathbf{\alpha} \right) = r_{1}T\left( \mathbf{v_{1}} \right) + r_{2}T\left( \mathbf{v_{2}} \right) + \cdots + r_{n}T\left( {\mathbf{v}}_{n} \right)\]

Tell me the image of the basis vectors without telling me the formula of
T, and I would still be able to compute the value of the output for any
given no input. If you have the image of basic for the T, then you can
compute the output for any given input.

\textbf{Property:} The linear transformation must fix a zero. This means
\(T\left( \mathbf{0} \right) = \mathbf{0}\) must be in the
codomain. This gives us a shortcut to check if a given function is
linear or not. If a zero gives an output as a nonzero vector, then you
can tell it is not a linear transformation.

Proof. Let
\(\mathbf{0} = \mathbf{v} - \mathbf{v}\text{.\ }\mathbf{v} \in \mathbb{R}^{n}\).

\[T\left( \mathbf{0} \right) = T\left( \mathbf{v} - \mathbf{v} \right) = T\left( \mathbf{v} \right) - T\left( \mathbf{v} \right) = \mathbf{0}\]

\(\Rightarrow ker(T)\) always contains the zero vector.


\section{Vector Spaces and Linear
Independence}

A homogeneous linear system can be represented as:

\[A\mathbf{x} = 0\]

Where \(A\) is an \(m \times n\) matrix, where
\(\mathbf{x} = \begin{bmatrix}
x_{1} \\
 \vdots \\
x_{n} \\
\end{bmatrix}\).

\[T:\mathbb{R}^{n} \rightarrow \mathbb{R}^{m}:T\left( \mathbf{x} \right) = A\mathbf{x}
\]

\textbf{Remark}

\begin{enumerate}
\def\labelenumi{\arabic{enumi}.}
\item
  Every homogeneous system is consistent (\(\geq 1\) soln)
\end{enumerate}

If the system gives you no solutions -- you know how that looks like.
But, for a homogeneous equation, the right side of a partitioned matrix
will always be zero for all components.

Moreover, the trivial solution always is a solution for
\(A\mathbf{x} = \mathbf{0}\). The reason why we call this the
trivial solution is that
\(A\left( \mathbf{0} \right) = \mathbf{0},\ \mathbf{x} = 0\)

\begin{enumerate}
\def\labelenumi{\arabic{enumi}.}
\setcounter{enumi}{1}
\item
  If \(A\sim H\) and \(H\) has a pivot in every column, then
  \(A\mathbf{x} = 0\) only has the trivial solution
  (\(\mathbf{x} = \mathbf{0})\).
\item
  If \(m < n\) (number of rows \textless{} columns) and \(A\sim H\)
  (\emph{H} is in RREF) then it is impossible for \(H\) to have a pivot
  in every column, which means a free variable definitely exist. This
  means \(A\mathbf{x} = \mathbf{0}\) has infinitely many
  solutions.
\end{enumerate}


\subsection{Null Space, row space, col
space}

\begin{definition}
Let \(A\) be an \(m \times n\) matrix

\begin{enumerate}
\def\labelenumi{\arabic{enumi}.}
\item
  The null space of \(A\) is a set of vectors \(\mathbf{x}\) such
  that \(A\mathbf{x} = 0\) i.e.,
  \(\left\{ \mathbf{x} \in \mathbb{R}^{n}\mid A\mathbf{x} = 0 \right\}\).
  It is the set of vectors that go to the zero vector when multiplying
  \(A\) from the right.
\item
  The row space of \(A\), \(row(A)\) is the span of row vectors of
  \(A\).
\item
  The column space of \(A\), \(col(A)\) is the span of column vectors of
  \(A\).
\end{enumerate}
\end{definition}

A linear system \(A\mathbf{x} = \mathbf{b}\) has a solution if
and only if \(\mathbf{b}\) is in the column space of \(A\).


\subsubsection{Deducing Null Spaces}

For example, suppose I have this matrix:

\[A = \begin{bmatrix}
2 & - 6 & - 2 & 4 \\
 - 1 & 3 & 3 & 2 \\
 - 1 & 3 & 7 & 10 \\
\end{bmatrix}\]

\begin{align*}
\text{row}(A) &= \text{span}\left( \lbrack 2,\  - 6,\ 2,\ 4\rbrack,\ \lbrack - 1,\ 3,\ 3,\ 2\rbrack,\ \lbrack - 1,\ 3,\ 7,\ 10\rbrack \right) \\
\text{col}(A) &= \text{span}\left( \begin{bmatrix}2 \\ - 1 \\ - 1 \\\end{bmatrix},\ \begin{bmatrix} - 6 \\3 \\3 \\\end{bmatrix},\ \begin{bmatrix} - 2 \\3 \\7 \\\end{bmatrix},\ \begin{bmatrix}4 \\2 \\10 \\\end{bmatrix} \right) \\
\text{null}(A) &= \left\{ \mathbf{x} \in \mathbb{R}^{4}\mid A\mathbf{x} = 0 \right\}
\end{align*}


We need to solve \(A\mathbf{x} = 0\).

\[\left\lbrack \begin{matrix}
2 & - 6 & - 2 & 4 \\
 - 1 & 3 & 3 & 2 \\
 - 1 & 3 & 7 & 10 \\
\end{matrix}\mid\,\begin{matrix}
0 \\
0 \\
0 \\
\end{matrix} \right\rbrack\sim\left\lbrack \begin{matrix}
1 & - 3 & 0 & 4 \\
0 & 0 & 1 & 2 \\
0 & 0 & 0 & 0 \\
\end{matrix}\;\middle|\;\begin{matrix}
0 \\
0 \\
0 \\
\end{matrix} \right\rbrack\]

There are infinitely many solutions to this.

\begin{align*}
x_{2} &= t \\
x_{4} &= s \\
x_{3} + 2x_{4} &= 0 \Rightarrow x_{3} = - 2x_{4} = - 2s \\
x_{1} - 3x_{2} + 4x_{4} &= 0 \Rightarrow x_{1} = 3x_{2} - 4x_{4} = 3t - 4s
 \\
\begin{bmatrix}x_{1} \\x_{2} \\x_{3} \\x_{4} \\\end{bmatrix} &= \begin{bmatrix}3t - 4s \\t \\ - 2s \\s \\\end{bmatrix} = \begin{bmatrix}3t \\t \\0 \\0 \\\end{bmatrix} + \begin{bmatrix} - 4s \\0 \\ - 2s \\s \\\end{bmatrix} \\
&= t\begin{bmatrix}3 \\1 \\0 \\0 \\\end{bmatrix} + s\begin{bmatrix} - 4 \\0 \\ - 2 \\1 \\\end{bmatrix}
\end{align*}


Therefore:

\[\text{null}(A) = \text{span}\left( \begin{bmatrix}
3 \\
1 \\
0 \\
0 \\
\end{bmatrix},\ \begin{bmatrix}
 - 4 \\
0 \\
 - 2 \\
1 \\
\end{bmatrix} \right)\]

Null space is the only space that requires computations.


\subsubsection{Proofs regarding null
spaces}

\textbf{Proposition:} Let \(A\mathbf{x} = \mathbf{b}\) be a
linear system with a particular solution \(p\).
\((A\mathbf{p} = \mathbf{b})\)

If \(\mathbf{h}\) is a vector in the \(null(A)\), then
\(\mathbf{p} + \mathbf{h}\) is also a solution for
\(A\mathbf{x} = \mathbf{b}\).

Proof.
\(A\left( \mathbf{p} + \mathbf{h} \right) = \mathbf{b}\)

\begin{align*}
\text{LHS} &= A\mathbf{p} + A\mathbf{h} \\
\text{LHS} &= A\mathbf{p} + 0 \\
\text{LHS} &= \mathbf{b}
\end{align*}


Remember that because h is a vector in the null space of A, it means
\(A\mathbf{h}\) is always equals to zero. \(A\mathbf{p} = b\) is
true because it is defined as true as part of the proof.

If \(\mathbf{q}\) is any solution to
\(A\mathbf{x} = \mathbf{b}\), then
\(\mathbf{q} = \mathbf{p} + {\mathbf{h}}_{1}\) for some
\(\mathbf{h} \in \text{null}(A)\)

\[A\mathbf{q} = A\left( \mathbf{p} + {\mathbf{h}}_{1} \right) = A\mathbf{p} + A{\mathbf{h}}_{1} = \mathbf{b} + 0 \Rightarrow {\mathbf{h}}_{1} \in \text{null}(A)\]


\subsection{Linear independence}

\textbf{Definition:} Linear independence (every single vector is useful;
none of the other vectors can be made up through linear combinations of
other vectors).

If
\({\mathbf{v}}_{1},\ {\mathbf{v}}_{2},\ \ldots,\ {\mathbf{v}}_{k} \in \mathbb{R}^{n}\),
\(r_{1}{\mathbf{v}}_{1} + r_{2}{\mathbf{v}}_{2} + \ldots + r_{k}{\mathbf{v}}_{k} = \mathbf{0}\)
has exactly 1 solution (that is \(\mathbf{r}\) being all zero), then
we say that \({\mathbf{v}}_{1},\ \ldots,\ {\mathbf{v}}_{k}\) are
linearly independent.

(It is impossible that this system has no solutions. If this system has
infinite solutions, then
\({\mathbf{v}}_{1},\ {\mathbf{v}}_{2},\ \ldots,\ {\mathbf{v}}_{k}\)
are linearly dependent.)


\subsubsection{Linear independence
properties}

\textbf{Remark:} \(r_{1} = r_{2} = r_{3} = \cdots = r_{k} = 0\) is a
solution for
\(r_{1}{\mathbf{v}}_{1} + r_{2}{\mathbf{v}}_{2} + \ldots + r_{k}{\mathbf{v}}_{k} = 0\),
and if that is the only solution then
\({\mathbf{v}}_{1},\ {\mathbf{v}}_{2},\ \ldots,\ {\mathbf{v}}_{k}\)
are linearly independent.

Suppose not. \(\alpha_{1},\ \alpha_{2},\ \ldots,\ \alpha_{k}\) is
another solution for
\(r_{1}{\mathbf{v}}_{1} + r_{2}{\mathbf{v}}_{2} + \ldots + r_{k}{\mathbf{v}}_{k} = 0\).
\(\exists\alpha_{i} \neq 0\) (Why are we writing that? Because one
solution has \(r_{1}\ldots r_{k}\) be all zero and we are trying to
incorporate a non-all-zero solution for
\(r_{1}{\mathbf{v}}_{1} + r_{2}{\mathbf{v}}_{2} + \ldots + r_{k}{\mathbf{v}}_{k} = 0\).
Our work is here:

\begin{align*}
\alpha_{1}{\mathbf{v}}_{1} + \alpha_{2}{\mathbf{v}}_{2} + \ldots + \alpha_{i}{\mathbf{v}}_{i} + \ldots + \alpha_{k}{\mathbf{v}}_{k} &= 0 \\
&\alpha_{1}{\mathbf{v}}_{1} + \alpha_{2}{\mathbf{v}}_{2} + \ldots  \\
&+ \alpha_{i - 1}{\mathbf{v}}_{i - 1} + \ldots + \alpha_{k}{\mathbf{v}}_{k}  \\
&= - \alpha_{i}{\mathbf{v}}_{i} \\
\frac{\alpha_{1}}{\alpha_{i}}{\mathbf{v}}_{1} + \frac{\alpha_{2}}{\alpha_{i}}{\mathbf{v}}_{2} + \cdots + \frac{\alpha_{k}}{\alpha_{i}}{\mathbf{v}}_{k} &= - {\mathbf{v}}_{i}
\end{align*}


If \({\mathbf{v}}_{i}\) can be represented by the left equation,
then \({\mathbf{v}}_{i}\) is redundant, which is \emph{impossible}
because we said that
\({\mathbf{v}}_{1},\ \ldots,\ {\mathbf{v}}_{k}\) are linearly
independent. We thus arrive at a contradiction.

\textbf{Definition:} Two vectors in \(\mathbb{R}^{n}\) are linearly
independent if they are non-zero and nonparallel.


\subsubsection{Testing for linear
independence}

\textbf{Example:}
\({\mathbf{v}}_{1} = \lbrack 2,\  - 1,\  - 1\rbrack,\ {\mathbf{v}}_{2} = \lbrack - 2,\ 3,\ 7\rbrack,\ {\mathbf{v}}_{3} = \lbrack 2,\ 1,\ 5\rbrack\)

Solution: If
\(r_{1}{\mathbf{v}}_{1} + r_{2}{\mathbf{v}}_{2} + r_{3}{\mathbf{v}}_{3} = 0\)
\textbf{has exactly one solution,} the vectors are linearly independent.
You can represent this in a matrix \(A\mathbf{r}\), where \(A\)'s
rows are formed by the column vectors \(\mathbf{v}\).

If a column is missing a pivot, then it is impossible for this system to
have a unique solution \(\mathbf{r}\).

\[\begin{bmatrix}
2 & - 2 & 2 \\
 - 1 & 3 & 1 \\
 - 1 & 7 & 5 \\
\end{bmatrix}\sim\begin{bmatrix}
1 & 0 & 0 \\
0 & 1 & 0 \\
2 & 1 & 0 \\
\end{bmatrix}\]

While we did not reduce the above matrix to RREF, but this is good
enough as we can see there is at least one column without the pivot.


\subsubsection{Testing for linear
dependence}

Because
\(r_{1}{\mathbf{v}}_{1} + r_{2}{\mathbf{v}}_{2} + r_{3}{\mathbf{v}}_{3} = 0\)
has more than one solution, the vectors are linearly dependent.

Is this matrix invertible (if no, then we've got linear dependence)?

\[\begin{bmatrix}
2 & - 1 & - 1 \\
 - 2 & 3 & 7 \\
2 & 1 & 5 \\
\end{bmatrix}\sim\begin{bmatrix}
1 & 0 & 2 \\
0 & 1 & 1 \\
0 & 0 & 0 \\
\end{bmatrix}\]

The rightmost column causes this matrix to have a free variable. If we
let it equal to zero:

\[\left\lbrack \begin{matrix}
1 & 0 & 2 \\
0 & 1 & 1 \\
0 & 0 & 0 \\
\end{matrix}\;\middle|\;\begin{matrix}
0 \\
0 \\
0 \\
\end{matrix} \right\rbrack\]

\(A\mathbf{x} = 0\) has the solution \(r\begin{bmatrix}
2 \\
1 \\
 - 1 \\
\end{bmatrix}\) (we factored out the negatives)

\[\left\lbrack {\mathbf{v}}_{3} \right\rbrack = 2\left\lbrack {\mathbf{v}}_{1} \right\rbrack + \left\lbrack {\mathbf{v}}_{2} \right\rbrack\]

\({\mathbf{v}}_{3}\) is a redundant vector, and so could the others?


\section{Subspaces and basis
vectors}

\emph{Hint: checking ``in a subspace'' WILL appear on an exam.}

\textbf{Definition:} Let \(W \subseteq \mathbb{R}^{n}\text{.\ W}\) is
called \textbf{a subspace} of \(\mathbb{R}^{n}\) if these following
properties hold:

\begin{enumerate}
\def\labelenumi{\arabic{enumi}.}
\item
  \(W \neq \varnothing\)
\item
  If \(\mathbf{u},\ \mathbf{v} \in W\), then
  \(\mathbf{u} + \mathbf{v} \in W\) (Add two vectors in the
  ``subspace'' and your sum must still be within the subspace. This is
  the closure property regarding subspaces.)
\item
  If \(\mathbf{u} \in W,\ r\mathbb{\in R}\), then
  \(r \cdot \mathbf{u} \in W\) (Closure under scalar multiplication)
\end{enumerate}

To check properties 2 and 3, you may need to pick variables arbitrarily
and you must end up with a way to represent the sum or scalar product of
the vectors in the way an individual vector in the set would be defined.

With all three holding \(W\) is a subspace.


\subsection{Checking if a set is a subspace:
Examples}

\textbf{Example}:
\(W = \left\{ \mathbf{v} \in \mathbb{R}^{2}\mid\mathbf{v} = \lbrack x,\ 0\rbrack,\ x\mathbb{\in R} \right\}\)

\emph{A subspace is usually seen as a line or an infinite plane, and so
on.}

Is \(W\) a subspace?

\begin{enumerate}
\def\labelenumi{\arabic{enumi}.}
\item
  \(W \neq \varnothing;\varnothing \subset W\) (NOT the zero vector.)
\item
  Pick \(\mathbf{u} = \left\lbrack x_{1},\ 0 \right\rbrack \in W\),
  \(\mathbf{v} = \left\lbrack x_{2},\ 0 \right\rbrack \in W\), The
  sum?
  \(\mathbf{u} + \mathbf{v} = \left\lbrack x_{1} + x_{2},\ 0 \right\rbrack \in W\)
  (it still lies on \(W\) and the \(x\)-axis.)
\item
  \(r\mathbf{u} = r\lbrack x,\ 0\rbrack = \lbrack rx,\ 0\rbrack \in W\)
\end{enumerate}

Because these properties hold, \(W\) is a subspace in
\(\mathbb{R}^{2}\).

\textbf{Second example:}
\(H = \{\mathbf{v} \in \mathbb{R}^{2} \mid \mathbf{v} = \lbrack 0,\ y\rbrack,\ y\mathbb{\in R\}}\)

You can show that \(H\) is a subspace in \(\mathbb{R}^{2}\).

\textbf{Third example:} \(H \cup W\)? No. The reason why this is not a
subspace.

We can disprove this with vector addition:

\[\lbrack 1,\ 0\rbrack + \lbrack 0,\ 1\rbrack = \lbrack 1,\ 1\rbrack \notin H \cup W\]

Therefore \(H \cup W\) can't be a subspace.


\textbf{Fourth example:} (Try this out of lecture) \(H \cap W\). That is
just the zero vector. Scalar multiply a zero vector by any amount, and
you will get the zero vector. Add the zero vector by another zero
vector, and you get the zero vector. Therefore, the zero subspace is
closed under addition and scalar multiplication.

\textbf{Remark:} I will check if \(\mathbf{0}\) is contained in
\(W\) \textbf{as a subspace must contain zero (but does not imply that a
subspace is valid).}

Suppose \(W\) is a subspace. This implies that \(W\) is non-empty,
implying \(\exists\mathbf{u} \in W\). This also implies:

\begin{align*}
&\Rightarrow - 1\mathbf{u} \in W \Rightarrow \mathbf{u} + \left( - 1\mathbf{u} \right) \in W \\
&\Rightarrow \mathbf{0} \in W
\end{align*}


If \(\overrightarrow{0}\) is not in the subspace, then I can immediately
declare that \(W\) is not a subspace. Because all subspaces must be
closed under addition and scalar multiplication you must be able to find
zero when doing a combination of the two.

\textbf{Example:}
\(W = sp\left\{ {\mathbf{w}}_{1},\ {\mathbf{w}}_{2},\ \ldots,\ {\mathbf{w}}_{k} \right\}\),
where
\(w_{1},\ w_{2},\ \ldots,\ w_{k} \in \mathbb{R}^{n} \Rightarrow W\) is a
subspace of \(\mathbb{R}^{n}\).

\begin{proof} Check if \(\mathbf{0}\) belongs to \(W\).

\[r_{1}{\mathbf{w}}_{1} + r_{2}{\mathbf{w}}_{2} + \ldots + r_{k}{\mathbf{w}}_{k} = \mathbf{0}\]

To show this is true, show that this system is homogeneous, and that
there is always one solution: \(r_{1} = r_{2} = \ldots = r_{k} = 0\).

Firstly, \(\mathbf{0}\) is a vector in \(W\), and a span means all
possible linear combinations. Therefore,
\(\mathbf{0} \in W,\  \Rightarrow \mathbf{w} \neq 0\).

Step 2: Check if \(\mathbf{u} + r\mathbf{v} \in W\) where
\(\mathbf{u} \cdot \mathbf{v} \in W,\ r\mathbb{\in R}\). If I
can answer this, then \(\mathbf{u},\ \mathbf{v}\) are closed
under vector addition and scalar multiplication.

\begin{align*}
\mathbf{u} &= \alpha_{1}{\mathbf{w}}_{1} + \alpha_{2}{\mathbf{w}}_{2} + \ldots + \alpha_{k}{\mathbf{w}}_{k} \\
\mathbf{v} &= \beta_{1}{\mathbf{w}}_{1} + \beta_{2}{\mathbf{w}}_{2} + \ldots + \beta_{k}{\mathbf{w}}_{k}
\end{align*}


Where
\(\alpha_{1}\mathbb{\in R\forall}i,\ \beta_{i}\mathbb{\in R\forall}i\)

\begin{align*}
&\mathbf{u} + r\mathbf{v}  \\
&= \left( \alpha{\mathbf{w}}_{1} + \ldots + \alpha_{k}{\mathbf{w}}_{k} \right)  \\
&+ \left( r\beta_{1}{\mathbf{w}}_{1} + r\beta_{2}{\mathbf{w}}_{2} + \ldots + r\beta_{k}{\mathbf{w}}_{k} \right) \\
& \\
&= \left( \alpha_{1} + r\beta_{1} \right){\mathbf{w}}_{1} + \left( \alpha_{2} + r\beta_{2} \right){\mathbf{w}}_{2}  \\
&+ \ldots + \left( \alpha_{n} + r\beta_{n} \right){\mathbf{w}}_{k}
\end{align*}


This is just another linear combination, showing that
\(\mathbf{u} + r\mathbf{v} \in sp\left\{ {\mathbf{w}}_{1}\ldots{\mathbf{w}}_{k} \right\} = W\)
which means \(\mathbf{u} + r\mathbf{v} \in W \Rightarrow\) \(W\)
is closed under addition and scalar multiplication. Therefore, \(W\) is
a subspace of \(\mathbb{R}^{n}\).



\end{proof}
\subsection{Finding a basis for a set of
vectors}

In a nutshell:

If we need to find a basis for
\(W = sp\left( w_{1},\ w_{2},\ \ldots,\ w_{k} \right)\)

\begin{enumerate}
\def\labelenumi{\arabic{enumi}.}
\item
  Form the matrix \(A\) whose \(j\)th column vector is \(w_{j}\)
\item
  Row-reduce \(A\) to REF, which we will call \(H\).
\item
  The set of all \(w_{j}\) such that the \(j\)th column of \(H\)
  contains a pivot is a basis for \(W\).
\end{enumerate}

\textbf{Definition:} Let \(W\) be a subspace of \(\mathbb{R}^{n}\).

If
\(Β = \left\{ {\mathbf{b}}_{1},\ {\mathbf{b}}_{2},\ \ldots,\ {\mathbf{b}}_{k} \right\}\)
is a subset of \(W\), then \(B\) is a \textbf{basis} for \(W\) if every
vector in \(W\) can be written \textbf{uniquely} as a linear combination
of vectors in \(B\). (All the vectors in \(W\) can be generated by
\(B\). If that holds, \(B\) has that property.)

\textbf{Remark:} Pick \(\mathbf{w} \in W\).

\[\mathbf{w} = r_{1}{\mathbf{b}}_{1} + r_{2}{\mathbf{b}}_{2} + \ldots + r_{k}{\mathbf{b}}_{k},\ r_{i}\mathbb{\in R\forall}i\]

And we need to prove this representation is unique, i.e.,
\(\{ r_{1},\ \ldots,\ r_{k}\}\) is unique. This gives us another linear
system, with unknowns \(\mathbf{r}\). We need to check if this
system has a unique solution; if it does, then those vectors
\({\mathbf{b}}_{1}\ldots{\mathbf{b}}_{k}\) are basis vectors.

\[\left\lbrack \begin{matrix}
{\mathbf{b}}_{1} & {\mathbf{b}}_{2} & \cdots & {\mathbf{b}}_{k} \\
\end{matrix}\;\middle|\;\begin{matrix}
\mathbf{0} \\
\end{matrix} \right\rbrack\ \]

This system having more than one solution indicates that
\(\{ r_{1}\ldots r_{k}\}\) is not unique.

\textbf{Example:}
\(W = sp\{{\mathbf{w}}_{1},\ {\mathbf{w}}_{2},\ \ldots,\ {\mathbf{w}}_{5}\}\)

\[w_{1} = \begin{bmatrix}
1 \\
0 \\
1 \\
0 \\
\end{bmatrix},\ w_{2} = \begin{bmatrix}
4 \\
1 \\
4 \\
1 \\
\end{bmatrix},\ w_{3} = \begin{bmatrix}
3 \\
1 \\
3 \\
2 \\
\end{bmatrix},\ w_{4} = \begin{bmatrix}
 - 1 \\
0 \\
 - 1 \\
1 \\
\end{bmatrix},w_{5} = \begin{bmatrix}
 - 6 \\
3 \\
 - 2 \\
0 \\
\end{bmatrix}\ \]

\[\begin{bmatrix}
1 & 0 & 0 & 0 & - 6 \\
0 & 1 & 0 & - 1 & 3 \\
0 & 0 & 1 & 1 & - 2 \\
0 & 0 & 0 & 0 & 0 \\
\end{bmatrix}\ \]

Then only
\({\mathbf{w}}_{1},\ {\mathbf{w}}_{2},\ {\mathbf{w}}_{3}\)
are single-pivot row vectors; all of the other vectors are redundant.
\({\mathbf{w}}_{4} = - {\mathbf{w}}_{2} + {\mathbf{w}}_{3}\)
and
\({\mathbf{w}}_{5} = - 6{\mathbf{w}}_{1} + {3\mathbf{h}}_{2} - 2{\mathbf{w}}_{3}\)

The basis vectors are
\(sp\left\{ {\mathbf{w}}_{1},\ {\mathbf{w}}_{2},\ {\mathbf{w}}_{3} \right\}\).

This is equivalent to
\(sp\left\{ {\mathbf{w}}_{1},\ {\mathbf{w}}_{2},\ {\mathbf{w}}_{3},\ {\mathbf{w}}_{4},\ {\mathbf{w}}_{5} \right\}\).
The span is in \(\mathbb{R}^{3}\);
\(\text{dimension}\left( sp\left\{ {\mathbf{w}}_{1},\ {\mathbf{w}}_{2},\ {\mathbf{w}}_{3} \right\} \right) = 3\).


\subsection{Basis Theorems}

\begin{theorem}
Let \(A \in M_{n \times n}\). The following are
equivalent (if you can show one of them all in this list will be true):

\begin{enumerate}
\def\labelenumi{\arabic{enumi}.}
\item
  The linear system \(A\mathbf{x} = \mathbf{b}\) has a unique
  solution for any \(\mathbf{b} \in \mathbb{R}^{n}\).
\item
  \(A\sim I\)
\item
  \(A\) is invertible.
\item
  The column vectors of \(A\) form a basis for \(\mathbb{R}^{n}\).
\item
  The row vectors in \(A\) form a basis for \(\mathbb{R}^{n}\).
\end{enumerate}
\end{theorem}


\subsubsection{Proofs for these
Theorems}

\begin{proof} \(1 \Rightarrow 2\) is the process of solving a linear
system. Reducing \(A\) to RREF, if every single column has a pivot
(\(n\) pivots), you'll get a unique solution. That matrix will look like
\(I\), thus forcing \(A\) to be row equivalent to the identity matrix.

This is trivially equivalent to \(3\).
\(\left\lbrack A\mid I \right\rbrack\sim\left\lbrack H\mid A^{- 1} \right\rbrack\),
where \(H\) is the identity matrix.

For no. 4: Column vectors form a basis \(\Rightarrow\) they are linearly
independent \(\Rightarrow\) every column must contain a pivot
\(\Rightarrow\) you will have \(n\) pivots, meaning no free variables
and a unique solution. If that is case, then \(A\sim I \Leftrightarrow\)
every single column has a pivot \(\Leftrightarrow\) they can form a
basis.

For no. 5: \(A\sim I \Leftrightarrow A^{T}\sim I\)



\end{proof}
\subsection{Rectangular cases}

\begin{theorem}
Let \(A\) be \(m \times n\) matrix. \((m > n)\). The
following are equivalent:

\begin{enumerate}
\def\labelenumi{\arabic{enumi}.}
\item
  Each \textbf{consistent} system \(A\mathbf{x} = \mathbf{b}\)
  has a unique solution.
\item
  The RREF of \(A\) must contain the \(n \times n\) identity matrix with
  \(m - n\) rows of zeros (on the bottom). This means there are
  \(m - n\) redundant equations.
\item
  The column vectors of \(A\) form a basis for the column space of
  \(A\).
\end{enumerate}
\end{theorem}


\subsubsection{Proofs for these
theorems}

\begin{proof} \(1 \Rightarrow 2\)

\(A\mathbf{x} = \mathbf{b}\), and for a unique solution, free
variables are not allowed to exist, which means every column must
contain a pivot. All entries below the identity matrix must be all zero,
implying \(m - n\) rows of zeroes.

\emph{Proof.} \(2 \Rightarrow 3\)

Every column has a pivot \(\Rightarrow\) the column vectors are linearly
independent \(\Rightarrow\) can serve as basis.

\emph{Proof.} \(3 \Rightarrow 1\)

Basis \(\Rightarrow\) linearly independent. Building a linear system out
of the basis guarantees no free variables \(\Rightarrow\) unique
solution.

\begin{remark}
The number of solutions will depend on whether free
variables exist. If you have a rectangular matrix \(M_{m < n}\) (wider
than it is tall), then \(A\mathbf{x} = \mathbf{b}\) is a linear
system with fewer equations than unknowns. If
\(A\mathbf{x} = \mathbf{b}\) is consistent (ruling out no
solution situations), then it has infinitely many solutions.
\end{remark}

\begin{theorem}
Let
\(B = \left\{ {\mathbf{b}}_{1},\ \ldots,\ {\mathbf{b}}_{k} \right\}\)
be a basis for the subspace \(W \subset \mathbb{R}^{n}\), if and only if
{[}1{]}
\(B = \left\{ {\mathbf{b}}_{1},\ \ldots,\ {\mathbf{b}}_{k} \right\}\)
is a linearly independent set and {[}2{]} it spans \(W\).
\end{theorem}

\emph{Proof.} \(B\) is a basis for \(W\) \(\Leftrightarrow\) Pick any
\(\mathbf{w} \in W\); \(\mathbf{w}\) has a unique
representation:
\(\mathbf{w} = r_{1}{\mathbf{b}}_{1} + r_{2}{\mathbf{b}}_{2} + \cdots + r_{k}{\mathbf{b}}_{k}\),
and the solution for \(r\) must be unique. You could view this as a
linear system \(A\mathbf{x} = \mathbf{b}\), that has a unique
solution where \(\mathbf{b} = \mathbf{w}\), and \(A\) is a
matrix with column vectors
\({\mathbf{b}}_{1}\ldots{\mathbf{b}}_{k}\), and
\(\mathbf{x} = \left\lbrack r_{1},\ r_{2},\ \ldots,\ r_{k} \right\rbrack\)
\(\Leftrightarrow\) no free variable \(\Leftrightarrow\) every column
must contain a pivot \(\Leftrightarrow\) none of any element in \(B\)
are redundant \(\Leftrightarrow\) \(\forall{\mathbf{b}}_{i}\)
(column vectors), it cannot be represented by linear combination of
other \({\mathbf{b}}_{i}\) \(\Leftrightarrow\) they are linearly
independent.



\end{proof}
\subsection{How to find a basis for any set of
vectors}

This: \(W = sp\{{\mathbf{w}}_{1}\cdots{\mathbf{w}}_{k}\}\) where
\({\mathbf{w}}_{i} \in \mathbb{R}^{n}\) for
\(i = 1,\ 2,\ \ldots,\ k\)

Step 1. Create a matrix \(A\) where the columns of it are formed by the
\(i\)th vectors in \(W\).

Step 2: \(A\sim H,\) where \(H\) is REF or RREF. (Idea: check if every
column has a pivot.

Step 3: The basis for \(W\) consists only of all the columns of \(A\),
corresponding to the columns of \(H\) that contains the pivots.

\textbf{Definition:} Let \(W\) be a subspace of \(\mathbb{R}^{n}\). The
number of elements in a basis for \(W\) is called the dimension of
\(W\), denoted by \(\dim(w)\). There are two facts supported by this
definition:

\textbf{{[}EXAM HINT{]} Fact 1:} Let \(W\) be a subspace of
\(\mathbb{R}^{n}\). Let
\({\mathbf{w}}_{1},\ {\mathbf{w}}_{2},\ \ldots,\ {\mathbf{w}}_{k}\)
be vectors in \(W\) that span \(W\) and let
\({\mathbf{v}}_{1},\ \ldots,\ {\mathbf{v}}_{m}\) be linearly
independent vectors in \(W\). Then, \(k \geq m\). (This fact is trivial
-- the dimension of \(W \geq K\).)

\({\mathbf{v}}_{1}\cdots{\mathbf{v}}_{m}\) being linearly
independent could be a basis for either \(W\) or a subspace of \(W\). In
other words,
\(sp\left\{ {\mathbf{v}}_{1}\cdots{\mathbf{v}}_{m} \right\} = U \leq W \Rightarrow m \leq k\).

TL;DR: The quantity of linear independent vectors in \(W\) is always
less or equal to the quantity of a list of vectors that have to at least
span \(W\).

\textbf{Fact 2:} Any two bases for a subspace \(W \in \mathbb{R}^{n}\)
contain the same number of vectors.

(The number of elements in a basis tells you the dimensions of the
subspace it is in. If you have two bases with different numbers of
vectors the dimensions formed by the span of it will be different.) For
example, for
\(\mathbb{R}^{2}:\ sp\left( \lbrack 1,0\rbrack,\ \lbrack 0,1\rbrack \right) = \mathbb{R}^{2}\).
Also,
\(sp\left( \lbrack 2,\ 0\rbrack,\ \lbrack 0,2\rbrack \right) = \mathbb{R}^{2}\),
as \(\begin{bmatrix}
2 & 0 \\
0 & 2 \\
\end{bmatrix}\begin{bmatrix}
x_{1} \\
x_{2} \\
\end{bmatrix} = \mathbf{b}\) has a unique solution, meaning every
\(\mathbf{b}\) can uniquely be represented by
\(\lbrack 2,\ 0\rbrack,\ \lbrack 0,\ 2\rbrack\).
\(sp\left( \lbrack 2,\ 0\rbrack,\ \lbrack 0,2\rbrack,\ \lbrack a,b\rbrack \right)\),
\(\lbrack a,b\rbrack = \frac{a}{2}\lbrack 2,\ 0\rbrack + \frac{b}{2}\lbrack 0,\ 2\rbrack \Rightarrow \lbrack a,\ b\rbrack\)
is a redundant vector.

\textbf{Proving Fact 1:} Given:

\begin{enumerate}
\def\labelenumi{\arabic{enumi}.}
\item
  \(sp\left\{ {\mathbf{w}}_{1}\cdots{\mathbf{w}}_{k} \right\} = W \in \mathbb{R}^{n}\).
\item
  \({\mathbf{v}}_{1}\cdots{\mathbf{v}}_{m} \in W\) are linearly
  independent
\end{enumerate}

Show \(k \geq m\).

Suppose \(k < m\). By the idea of span,
\(\forall{\mathbf{v}}_{i}\ \left( i \in \lbrack 1,\ 2,\ 3,\ \ldots,\ m\rbrack \right) \in W = sp\left\{ {\mathbf{w}}_{1}\cdots{\mathbf{w}}_{k} \right\}\).
All of \(v_{i}\) are chosen from \(W\). This means

\begin{align*}
{\mathbf{v}}_{1} &= a_{11}{\mathbf{w}}_{1} + a_{21}{\mathbf{w}}_{2} + \cdots + a_{k1}{\mathbf{w}}_{k} \\
{\mathbf{v}}_{2} &= a_{12}{\mathbf{w}}_{1} + a_{22}{\mathbf{w}}_{2} + \cdots + a_{k2}{\mathbf{w}}_{k} \\
&\vdots  \\
{\mathbf{v}}_{m} &= a_{1m}{\mathbf{w}}_{1} + a_{2m}{\mathbf{w}}_{2} + \cdots + a_{km}{\mathbf{w}}_{k}
\end{align*}


This is true by definition of span (all in \(\mathbf{v}\) are in the
span of \(W\)).

Pick a vector \(\mathbf{w} \in W\).

\[{\mathbf{w} }{= r_{1}{\mathbf{v}}_{1} + r_{2}{\mathbf{v}}_{2} + \cdots + r_{m}{\mathbf{v}}_{m} }{\in sp\left\{ {\mathbf{v}}_{1}\cdots{\mathbf{v}}_{m} \right\}} \]

By the definition of span. Now, we can rewrite \(\mathbf{w}\) by
replacing \({\mathbf{v}}_{1}\) with the first equation, and so on.

\[{\mathbf{w} }{= r_{1}\left( a_{11}{\mathbf{w}}_{1} + a_{21}{\mathbf{w}}_{2} + \cdots a_{k1}{\mathbf{w}}_{k} \right) }{+ r_{2}\left( a_{12}{\mathbf{w}}_{1} + a_{22}{\mathbf{w}}_{2} + \cdots + a_{k2}{\mathbf{w}}_{k} \right) }{+ \cdots + r_{n}(a_{1m}{\mathbf{w}}_{1} + a_{2m}{\mathbf{w}}_{2} }{+ \cdots + a_{km}{\mathbf{w}}_{k})} \]

\[{}{= \left( r_{1}a_{11} + r_{2}a_{12} + \cdots + r_{m}a_{1m} \right){\mathbf{w}}_{1} }{+ \cdots + \left( r_{1}a_{k1} + r_{2}a_{k2} + \cdots + r_{m}a_{km} \right){\mathbf{w}}_{k}} \]

If \(\mathbf{w} = 0\), then \(A\mathbf{r} = \mathbf{0}\),
where \(A\) is a matrix formed by the column vectors of
\({\mathbf{v}}_{1}\cdots{\mathbf{v}}_{m}\). We know that
\({\mathbf{v}}_{1}\cdots{\mathbf{v}}_{m}\) is linearly
independent, meaning the matrix has a unique and trivial solution. This
means \(\mathbf{r} = 0\). This means

\begin{align*}
r_{1}a_{11} + r_{2}a_{12} + \cdots + r_{m1}a_{1m} &= 0 \\
&\vdots  \\
r_{1}a_{k1} + r_{2}a_{k2} + \cdots + r_{m}a_{km} &= 0
\end{align*}


This system only contains one solution, and that solution must be the
trivial solution of \(\mathbf{r} = 0\). However, you have \(k\)
equations and \(m\) unknowns. Recall that we assumed \(k < m\), meaning
the number of equations is less than the number of unknowns, and you
cannot have a unique solution that way -- you would have infinitely many
solutions.

\sout{You are supposed to have infinitely many} \(\mathbf{r}\)
\sout{to satisfy}
\(r_{1}{\mathbf{v}}_{1} + r_{2}{\mathbf{v}}_{2} + \cdots + r_{m}{\mathbf{v}}_{m} = \mathbf{0}\)\sout{.
Therefore} \({\mathbf{v}}_{1}\cdots{\mathbf{v}}_{m}\) \sout{are
linearly independent.}

\textbf{{[}1{]} Remark.} (Existence and determination of Base): Every
subspace \(W \subset \mathbb{R}^{n}\) has a basis and
\(\dim(w) \leq n\).

Pick \(\mathbf{b} \in W\), and pick a collection of linearly
independent vectors
\({\mathbf{w}}_{1},\ {\mathbf{w}}_{2},\ \ldots,\ {\mathbf{w}}_{k}\),
and use them to build a linear system. The objective is to make this
system always have a solution.

\[{\mathbf{w}}_{1}x_{1} + {\mathbf{w}}_{2}x_{2} + ⋯ + {\mathbf{w}}_{k}x_{k} = \mathbf{b}\]

Should have a unique solution \(\Rightarrow\) pivot for every single
column vector. If so, then they form a basis.

\textbf{{[}2{]}} Every linearly independent set of vectors in
\(\mathbb{R}^{n}\) can be enlarged if necessary to become a basis for
\(\mathbb{R}^{n}\).

\(sp\left( {\mathbf{v}}_{1}\cdots{\mathbf{v}}_{k} \right)\)
gives you a \(k\)-dimensional space, if \(k < n\). If you add some more
linearly independent vectors (to all of the existing span) so that the
number of vectors matches \(n\), i.e.

\[sp\left( {\mathbf{v}}_{1},\ \ldots,\ {\mathbf{v}}_{k},\ {\mathbf{v}}_{k + 1},\ \ldots,\ {\mathbf{v}}_{n} \right)\]

This forms a basis for \(\mathbb{R}^{n}.\)

\emph{Recall: Number of basis vectors of a subset gives you, its
dimension.}

\textbf{{[}3{]}} If \(W\) is a subspace of
\(\mathbb{R}^{n},\dim(w) = k\). Then:

\begin{enumerate}
\def\labelenumi{\alph{enumi})}
\item
  Every linearly independent \(k\) vectors in \(W\) is a basis for \(W\)
\item
  Every \(k\) vectors, the span of them gives you \(W\), then those
  \(k\) vectors are basis for \(W\).
\end{enumerate}


\subsection{Finding a basis from a
matrix}

\[A = \begin{bmatrix}
1 & 1 & 2 & 1 & 2 \\
2 & - 1 & 1 & 1 & 1 \\
1 & 2 & 3 & 1 & 3 \\
2 & 1 & 3 & 3 & 3 \\
\end{bmatrix}\]

What is the basis for the row space, and what is the basis for the
column space?

The candidate bases for
\(\text{col}(A) = sp\left\{ {\mathbf{c}}_{1},\ {\mathbf{c}}_{2},\ {\mathbf{c}}_{3},\ {\mathbf{c}}_{4},\ {\mathbf{c}}_{5} \right\}\),
\(\text{row}(A) = sp\left\{ {\mathbf{R}}_{1},\ {\mathbf{R}}_{2},\ {\mathbf{R}}_{3},\ {\mathbf{R}}_{4} \right\}\)

To figure out the basis, we need to figure which \(\mathbf{c}\) and
which \(\mathbf{r}\) is redundant. Row reducing this matrix:

\[A\sim H = \begin{bmatrix}
1 & 0 & 1 & 0 & 1 \\
0 & 1 & 1 & 0 & 1 \\
0 & 0 & 0 & 1 & 0 \\
0 & 0 & 0 & 0 & 0 \\
\end{bmatrix}\]

Columns 1, 2, and 4 have pivots. Columns 3 and 5 are redundant and can
be made up by linear combinations of columns 1, 2, and 4. For example:

\begin{align*}
{\mathbf{c}}_{5} &= {\mathbf{c}}_{1} + {\mathbf{c}}_{2} \\
{\mathbf{c}}_{3} &= {\mathbf{c}}_{1} + {\mathbf{c}}_{2}
\end{align*}


We can remove the redundant vectors from \(\text{col}(A)\) to determine
the basis for the column vectors of \(A\):

\[\text{col}(A) = sp\left\{ {\mathbf{c}}_{1},\ {\mathbf{c}}_{2},\ {\mathbf{c}}_{4} \right\} \Rightarrow \dim\left( \text{col}(A) \right) = 3\]

The row space of \(A\): Remove \({\mathbf{R}}_{4}\); it is
pivotless.

\[\text{row}(A) = sp\left\{ {\mathbf{R}}_{1},\ {\mathbf{R}}_{2},\ {\mathbf{R}}_{3} \right\} \Rightarrow \dim\left( \text{row}(A) \right) = 3\]


\subsubsection{The general procedure}

How to find bases for spaces associated with a matrix?

Let \(A \in M_{m \times n}\) matrix, and \(A\sim H,\ H\) is RREF of
\(A\).

There are 3 types of spaces associated with \(A\): \(A\) has

\begin{enumerate}
\def\labelenumi{\alph{enumi})}
\item
  Row space
\item
  Column space
\item
  Null space (span of vectors for \(\mathbf{x}\) that satisfies
  \(A\mathbf{x} = 0\))
\end{enumerate}

Procedure:

\begin{enumerate}
\def\labelenumi{\alph{enumi})}
\item
  The non-zero rows of \(H\) form a basis for the row space.
\item
  The column has a pivot \(\Rightarrow\) the column is a basis vector
  for the column space.
\item
  For a basis of null space \(\mathcal{N}\), we need to solve this
  system: \(H\mathbf{x} = 0\), to see how many basis vectors are in
  the solution space. The quantity of free variables in the system above
  will determine how many basis vectors in the solution space. If there
  are no free variables, then you only have one vector as the null
  space. One free variable generates a line; 2 generates a plane, 3
  generates \(\mathbb{R}^{3}\), and so on.
\end{enumerate}

\emph{Recall: Number of basis vectors of a subset gives you, its
dimension.}

\[\text{no.\ of\ basis\ for\ N} = \dim\left( \mathcal{N} \right) = \text{no.\ of\ free\ variables} = \text{no.\ of\ pivot-free\ columns}\]

\begin{remark}
\[\dim\left( \text{col}(A) \right) = \dim\left( \text{row}(A) \right) = \text{no.\ of\ pivot\ columns\ in\ H}\]
\end{remark}

\begin{example}
\[A = \begin{bmatrix}
1 & 1 & 2 & 1 & 2 \\
2 & - 1 & 1 & 1 & 1 \\
1 & 2 & 3 & 1 & 3 \\
2 & 1 & 3 & 3 & 3 \\
\end{bmatrix}\sim H = \begin{bmatrix}
1 & 0 & 1 & 0 & 1 \\
0 & 1 & 1 & 0 & 1 \\
0 & 0 & 0 & 1 & 0 \\
0 & 0 & 0 & 0 & 0 \\
\end{bmatrix}\]
\end{example}

Compute \(\text{null}(A)\). Firstly, column 3 and column 5 of \(H\) are
free variables, so we have two free variables. Therefore, the dimension
of the null space is 2.

\[\begin{bmatrix}
1 & 0 & 1 & 0 & 1 \\
0 & 1 & 1 & 0 & 1 \\
0 & 0 & 0 & 1 & 0 \\
0 & 0 & 0 & 0 & 0 \\
\end{bmatrix}\begin{bmatrix}
x_{1} \\
x_{2} \\
x_{3} \\
x_{4} \\
x_{5} \\
\end{bmatrix} = \begin{bmatrix}
0 \\
0 \\
0 \\
0 \\
\end{bmatrix} \Leftrightarrow \ \left\lbrack \begin{matrix}
1 & 0 & 1 & 0 & 1 \\
0 & 1 & 1 & 0 & 1 \\
0 & 0 & 0 & 1 & 0 \\
0 & 0 & 0 & 0 & 0 \\
\end{matrix}\;\middle|\;\begin{matrix}
0 \\
0 \\
0 \\
0 \\
\end{matrix} \right\rbrack\]

Solve this linear system:

\begin{align*}
1 \cdot x_{4} &= 0 \Rightarrow x_{4} = 0 \\
x_{3} &= t \\
x_{5} &= s \\
x_{1} &= - x_{3} - x_{5} = - t - s \\
x_{2} &= - x_{3} - x_{5} = - t - s \\
\begin{bmatrix}x_{1} \\x_{2} \\x_{3} \\x_{4} \\x_{5} \\\end{bmatrix} &= \begin{bmatrix} - t - s \\ - t - s \\t \\0 \\s \\\end{bmatrix} = t\begin{bmatrix} - 1 \\ - 1 \\1 \\0 \\0 \\\end{bmatrix} + s\begin{bmatrix} - 1 \\ - 1 \\0 \\0 \\1 \\\end{bmatrix}
\end{align*}


Therefore, \(\text{Null}(A) = sp\left\{ \begin{bmatrix}
 - 1 \\
 - 1 \\
1 \\
0 \\
0 \\
\end{bmatrix},\begin{bmatrix}
 - 1 \\
 - 1 \\
0 \\
0 \\
1 \\
\end{bmatrix} \right\}\)

The rank of a matrix is how many basis vectors it can form. Also, it can
be quantified by the number of pivots in the matrix when reduced to
RREF.


\subsubsection{Rank}

Let \(A \in M_{m \times n}\) matrix.

\begin{enumerate}
\def\labelenumi{\arabic{enumi}.}
\item
  The dimension of \(\text{col}(A) = \dim\left( \text{row}(A) \right)\)
  is called the rank of \(A\), denoted by \(\text{rank}(A)\). The rank
  can also be seen as the pivot count for a matrix reduced to RREF.
\item
  The dimension of the nullspace of \(A\), represented as
  \(\text{null}(A)\) is called the nullity of \(A\).
\end{enumerate}

\begin{theorem}
The rank equation of a matrix, where \(H\) is in REF
or RREF:

\begin{align*}
&A\sim H \\
&A,\ H \in M_{m \times n}\ \left( \mathbb{R} \right)
\end{align*}

\end{theorem}

Based on \(H\), we have:

\[\text{rank}(A) + \text{nullity}(A) = n\ \text{(width\ of\ matrix)}\]

\begin{proof} If \(A\sim H\)

\[\text{rank}(A) = \text{no.\ of\ pivots\ in\ H} = \text{no.\ of\ columns\ in\ H\ with\ pivots}\]

\[\text{nullity}(A) = \text{no.\ of\ pivotless\ columns}\]

\begin{align*}
\text{rank}(A) + \text{nullity}(A) &= \text{col\ with\ pivots\ in\ H} + \text{col\ without\ pivots\ in\ H} \\
&= \text{no.\ of\ columns\ in\ A\ and\ H} \\
&= n
\end{align*}


If you have a matrix with a zero column on the right:

\[\left\lbrack \begin{matrix}
1 & 2 & 3 & 0 \\
4 & 5 & 6 & 0 \\
7 & 8 & 9 & 0 \\
10 & 11 & 12 & 0 \\
\end{matrix}\;\middle|\;\begin{matrix}
0 \\
0 \\
0 \\
a \\
\end{matrix} \right\rbrack = \mathbf{0}\]

\(a\) is guaranteed to be mapped to zero. That is, the last element of
\(\mathbf{x}\).

\[x \cdot a = b \Rightarrow x = \frac{b}{a};b \neq 0 \Rightarrow a \neq 0\]



\end{proof}
\section{Matrices and Linear Maps}

Recall from section 1: A linear map \(T\) from \(\mathbb{R}^{n}\) to
\(\mathbb{R}^{m}\) is a map that meets the following conditions:

\[T\left( \mathbf{x} + \alpha\mathbf{y} \right) = T\left( \mathbf{x} \right) + \alpha T(\mathbf{y})\]

As long as the map satisfies this then it is a linear map. This map can
be modeled with a matrix of the appropriate size.

If \(T:\mathbb{R}^{n} \rightarrow \mathbb{R}^{m}\) is a linear map, let
\(A \in M_{m \times n}\ \left( \mathbb{R} \right)\).

\[A = \begin{bmatrix}
 \mid & \mid & \mid & \cdots & \mid \\
T({\mathbf{e}}_{1}) & T({\mathbf{e}}_{2}) & T({\mathbf{e}}_{3}) & \cdots & T({\mathbf{e}}_{n}) \\
 \mid & \mid & \mid & \cdots & \mid \\
\end{bmatrix}\]

\(T\left( {\mathbf{e}}_{i} \right) = {\mathbf{v}}_{i} \in \mathbb{R}^{m};{\mathbf{v}}_{i}\)
has \(m\) rows.
\({\mathbf{e}}_{1},\ {\mathbf{e}}_{2},\ \ldots,\ {\mathbf{e}}_{n}\)
are basis vectors that span \(\mathbb{R}^{n}\). (i.e.,
\(sp\left( {\mathbf{e}}_{1},\ {\mathbf{e}}_{2},\ \ldots,\ {\mathbf{e}}_{n} \right) = \mathbb{R}^{n}\)).

Then \(T\left( \mathbf{x} \right) = A\mathbf{x}\), where
\(\mathbf{x} \in \mathbb{R}^{n}\).

You can call \(A\) as the standard matrix representation of \(T\).

\emph{Proof:} Why do linear maps work?

For any possible input
\(\mathbf{x} \in \mathbb{R}^{n},\ \exists!\mathbf{x} = x_{1}e_{1} + x_{2}e_{2} + \cdots + x_{n}{\mathbf{e}}_{n}\)
(\(\mathbf{x}\) can be represented \emph{uniquely} in the fashion of
a linear combination, by definition of basis.)

\begin{align*}
T\left( \mathbf{x} \right) &= T\left( x_{1}{\mathbf{e}}_{1} + x_{2}{\mathbf{e}}_{2} + \cdots + x_{n}{\mathbf{e}}_{n} \right) \\
&= x_{1}T\left( {\mathbf{e}}_{1} \right) + x_{2}T\left( {\mathbf{e}}_{2} \right) + \cdots + x_{n}T\left( {\mathbf{e}}_{n} \right) \\
&= \begin{bmatrix} \mid & \mid & \mid & \cdots & \mid \\T\left( {\mathbf{e}}_{1} \right) & T\left( {\mathbf{e}}_{2} \right) & T\left( {\mathbf{e}}_{3} \right) & \cdots & T\left( {\mathbf{e}}_{n} \right) \\ \mid & \mid & \mid & \cdots & \mid \\\end{bmatrix}\begin{bmatrix}x_{1} \\x_{2} \\ \vdots \\x_{n} \\\end{bmatrix} \\
&= x_{1}\begin{bmatrix} \mid \\T\left( {\mathbf{e}}_{1} \right) \\ \mid \\\end{bmatrix} + x_{2}\begin{bmatrix} \mid \\T\left( {\mathbf{e}}_{2} \right) \\ \mid \\\end{bmatrix} + \cdots + x_{n}\begin{bmatrix} \mid \\T\left( {\mathbf{e}}_{n} \right) \\ \mid \\\end{bmatrix}
\end{align*}


Example: \(T:\mathbb{R}^{4} \rightarrow \mathbb{R}^{3}\) is a linear
map.

\[{T:\left( \left\lbrack x_{1},\ x_{2},\ x_{3},\ x_{4} \right\rbrack \right) }{= \left\lbrack x_{2} - 3x_{3} - x_{4},\ 6x_{1} + 5x_{2},\ x_{3} + 2x_{1} \right\rbrack} \]

How do I redefine \(T\) in terms of the matrix \(A\) such that
\(T\left( \mathbf{x} \right) = A\mathbf{x}\), i.e.

\[T\left( \left\lbrack x_{1},\ x_{2},\ x_{3},\ x_{4} \right\rbrack \right) = A \cdot \begin{bmatrix}
x_{1} \\
x_{2} \\
x_{3} \\
x_{4} \\
\end{bmatrix}\]

Solution: \(\text{domain}(T) = \mathbb{R}^{4}\).
\(\mathbb{R}^{4} = sp\left( \lbrack 1,\ 0,\ 0,\ 0\rbrack,\ \lbrack 0,\ 1,\ 0,\ 0\rbrack,\ \lbrack 0,\ 0,\ 1,\ 0\rbrack,\ \lbrack 0,\ 0,\ 0,\ 1\rbrack \right)\)

To construct \(A\), compute the image of the basis:
\(T({\mathbf{e}}_{1\ to\ 4})\):

\begin{align*}
T\left( {\mathbf{e}}_{1} \right) &= \lbrack 0,\ 6,\ 2\rbrack \\
T\left( {\mathbf{e}}_{2} \right) &= \lbrack 1,\ 5,\ 0\rbrack \\
T\left( {\mathbf{e}}_{3} \right) &= \lbrack - 3,\ 0,\ 1\rbrack \\
T\left( {\mathbf{e}}_{4} \right) &= \lbrack - 1,\ 0,\ 0\rbrack
\end{align*}


(\({\mathbf{e}}_{1} = \lbrack 1,\ 0,\ 0,\ 0\rbrack\), and so on --
we just ran those through the linear map.)

\[{A }{= \begin{bmatrix}
 \mid & \mid & \mid & \mid \\
T\left( {\mathbf{e}}_{1} \right) & T\left( {\mathbf{e}}_{2} \right) & T\left( {\mathbf{e}}_{3} \right) & T\left( {\mathbf{e}}_{4} \right) \\
 \mid & \mid & \mid & \mid \\
\end{bmatrix} }{= \begin{bmatrix}
0 & 1 & - 3 & 1 \\
6 & 5 & 0 & 0 \\
2 & 0 & 1 & 0 \\
\end{bmatrix}_{\text{this\ is\ A}}\begin{bmatrix}
x_{1} \\
x_{2} \\
x_{3} \\
x_{4} \\
\end{bmatrix}} \]

\[\Rightarrow T\left( \mathbf{x} \right) = A\mathbf{x}\]

Question: If we have a linear map
\(T:\mathbb{R}^{n} \rightarrow \mathbb{R}^{m}\), how can the
\(A(T\left( \mathbf{x} \right) = A\mathbf{x})\) help you to find
\(\ker(T)\) and \(\text{range}{(T)}\)?

Answer: \(\ker(T) = \text{null}{(A)}\), and
\(\text{range}{(T)} = \text{col}(A)\)

Why?

\[\ker(T) = \left\{ \mathbf{x} \in \mathbb{R}^{n}\mid T\left( \mathbf{x} \right) = 0 \right\} = \left\{ \mathbf{x} \in \mathbb{R}^{n}\mid A\mathbf{x} = 0 \right\} = \text{null}(A)\]

You can tell \(\text{null}(A)\) by checking
\(\left( A\sim H_{\text{RREF}} \right)\) the location of the pivots and
the number of the pivots in \(H_{\text{RREF}}\).

\[\text{range}{(T)} = \left\{ T\left( \mathbf{x} \right) \in \mathbb{R}^{M}\mid\mathbf{x} \in \mathbb{R}^{n} \right\} = \left\{ A\mathbf{x}\mid\mathbf{x} \in \mathbb{R}^{n} \right\} = \text{col}(A)\]

(How many columns have a pivot?)

Example: The linear map
\(T:\left( \left\lbrack x_{1},\ x_{2},\ x_{3},\ x_{4} \right\rbrack \right) = \left\lbrack x_{2} - 3x_{3} - x_{4},\ 6x_{1} + 5x_{2},\ x_{3} + 2x_{1} \right\rbrack\):
Deduce the kernel and the range. You can do that by finding the basis.

Firstly, find matrix \(A\) (already deduced before). Reduce it to \(H\).

\[\begin{bmatrix}
0 & 1 & - 3 & 1 \\
6 & 5 & 0 & 0 \\
2 & 0 & 1 & 0 \\
\end{bmatrix}\sim\begin{bmatrix}
1 & 0 & 0 & - 5/24 \\
0 & 1 & 0 & 1/4 \\
0 & 0 & 1 & 5/12 \\
\end{bmatrix} = H\]

\[T\left( \mathbf{x} \right) = A\mathbf{x}\]

\(\mathbf{c}\) represents the column vectors of \(H\).

\begin{align*}
{\mathbf{c}}_{4} &= - \frac{5}{24}{\mathbf{c}}_{1} + \frac{1}{4}{\mathbf{c}}_{2} + \frac{5}{12}{\mathbf{c}}_{3} \\
{\mathbf{c}}_{4} + \frac{5}{24}{\mathbf{c}}_{1} - \frac{1}{4}{\mathbf{c}}_{2} - \frac{5}{12}{\mathbf{c}}_{3} &= 0
\end{align*}


\[\text{col}(A) = \text{span}\left( c_{1},\ c_{2},\ c_{3} \right) = \text{range}(T)\]

\[\ker(T) = \left\{ \mathbf{x}\mid T\left( \mathbf{x} \right) = 0 \right\} = \left\{ \mathbf{x}\mid A\mathbf{x} = 0 \right\} \Leftrightarrow \left\{ \mathbf{x}\mid H\mathbf{x} = 0 \right\}\]

The \(\Leftrightarrow\) statement is true as
\(\text{null}{(A)} = \text{null}(H)\) if \(A\sim H\).

To figure out the kernel, I need to figure out what is happening to
\(H\mathbf{x} = 0\).

\[H\mathbf{x} = 0\]

\[\begin{bmatrix}
1 & 0 & 0 & - 5/24 \\
0 & 1 & 0 & 1/4 \\
0 & 0 & 1 & 5/12 \\
\end{bmatrix}\begin{bmatrix}
x_{1} \\
x_{2} \\
x_{3} \\
x_{4} \\
\end{bmatrix} = \begin{bmatrix}
0 \\
0 \\
0 \\
\end{bmatrix}\]

\begin{align*}
x_{1} - \frac{5}{24}x_{4} &= 0 \Rightarrow x_{1} = \frac{5}{24}x_{4} = \frac{5}{24}s \\
x_{2} + \frac{1}{4}x - 4 &= 0 \Rightarrow x_{2} = - \frac{1}{4}x_{4} = - \frac{1}{4}s \\
x_{3} + \frac{5}{12}x_{4} &= 0 \Rightarrow x_{3} = - \frac{5}{12}x_{4} = - \frac{5}{12}s \\
x_{4}\ \text{is\ a\ free\ variable},\ x_{4} &= s
\end{align*}


\[\begin{bmatrix}
x_{1} \\
x_{2} \\
x_{3} \\
x_{4} \\
\end{bmatrix} = \left\{ s\begin{bmatrix}
5/24 \\
 - 1/4 \\
 - 5/12 \\
1 \\
\end{bmatrix}\mid s\mathbb{\in R} \right\} = sp\left\{ \begin{bmatrix}
5/24 \\
 - 1/4 \\
 - 5/12 \\
1 \\
\end{bmatrix} \right\} = \ker(T)\]

This is the basis for the kernel, or a basis for
\(\text{null}{(A)} = \text{null}(H)\). \(\text{nullity}{(A)} = 1\)

\[\text{nullity}(A) + \text{rank}(A) = 1 + 3 = 4 = \text{no.\ of\ columns\ in\ A}\]

Based on this observation:

\emph{Remark.} Let \(T:\mathbb{R}^{n} \rightarrow \mathbb{R}^{m}\) be a
linear map.

\[T\left( \mathbf{x} \right) = A\mathbf{x}\]

We can define

\begin{enumerate}
\def\labelenumi{\arabic{enumi}.}
\item
  \(\text{rank}(T)\) is the dimension of the range of \(T\), which is
  the same as \(\text{rank}(A)\).
\item
  The nullity of \(T\) (\(\text{nullity}{(T)})\) is the dimension of the
  kernel, or the same as \(\text{nullity}{(A)}\).
\item
  If \(A\) is an invertible matrix, then \(T\) is invertible as well.
  (i.e., \(\exists T^{- 1})\).
  \(T^{- 1}\left( \mathbf{x} \right) = A^{- 1}\mathbf{x}\)
\end{enumerate}

Using these definitions, we can determine the fundamental theorem of the
linear map:

\begin{theorem}
\(\text{rank}(T) + \text{nullity}{(T)} = \dim\left( \text{domain}(T) \right)\)
\end{theorem}

This is the same as
\(\text{rank}(A) + \text{nullity}{(A)} = \text{no.\ of\ columns\ in\ A}\)

We have a new way to say matrix \(A\) is invertible:

\begin{align*}
&\left( A \in M_{n \times n} \right) \\
\text{rank}(A) &= n \\
A\mathbf{x} &= \mathbf{b}\ \text{has\ a\ unique\ sol'n} \Rightarrow \text{no\ free\ variable} \Rightarrow n\ \text{pivots} \\
\Leftrightarrow \text{nullity}(A) &= 0 \\
\Leftrightarrow A\mathbf{0} &= \mathbf{0} \\
\Leftrightarrow \text{rank}(A) &= \text{no.\ of\ columns}
\end{align*}


In this case, we call \(A\) full rank. The full rank is the highest
possible value for the rank.

\(A\) is invertible if and only if \(A\) is full rank.

If a space only contains the zero vector, the dimension is zero.

\begin{theorem}
Preservation of subspaces.
\end{theorem}

Let \(T:\mathbb{R}^{n} \rightarrow \mathbb{R}^{m}\) be a linear map.

If \(W\) is a subspace of \(\mathbb{R}^{n}\), then \(T(W)\) is a
subspace in \(\mathbb{R}^{n}\). (Linear transformations never knock
inputs within a subspace out of the subspace.)

Proof. \(W = \left\{ \begin{matrix}
 \neq \varnothing \\
\mathbf{x},\ \mathbf{y} \in W,\ \mathbf{x} + \alpha\mathbf{y} \in W \\
\end{matrix} \right.\ \)

How about \(T(W)\)? If
\(\mathbf{\alpha},\ \mathbf{\beta} \in T(W)\), we want to check
\(\alpha + r\beta \in T(W)\), and \(r\mathbb{\in R}\)

\[\exists\mathbf{x},\ \mathbf{y} \in W,\ T\left( \mathbf{x} \right) = \mathbf{\alpha},\ T\left( \mathbf{y} \right) = \mathbf{\beta}\]

This means you have two inputs in the domain of \(T\) such that (the
above statement).

We can find the inputs corresponding to
\(\mathbf{\alpha} + r\mathbf{\beta}\)

Consider \(\mathbf{x} + r\mathbf{y} \in W\), as \(W\) is a
subspace.

\[T\left( \mathbf{x} + r\mathbf{y} \right) = T\left( \mathbf{x} \right) + rT\left( \mathbf{y} \right) = \mathbf{\alpha} + r\mathbf{\beta} \in T(w)\]

This means \(T(w)\) is closed under addition and scalar multiplication.

\(T(W) \neq \varnothing\)\emph{,} as \(W\) is a subspace
\(\Rightarrow \mathbf{0} \in W\) (which again shows
\(W \neq \varnothing\))

\[T\left( \mathbf{0} \right) = \mathbf{0} \Rightarrow T(\mathbf{w}) \neq \varnothing\]

\(\mathbf{0}\) is an input in W.

Homework (something you should be able to do for the final exam): Prove
if \(W'\) is a subspace in \(\mathbb{R}^{m}\), then the pre-image of
\(W'\) under \(T\) is a subspace of the domain \(\mathbb{R}^{n}\).

\(f:V \rightarrow W\): if \(f\) is bijective, then \(|V| = |W|\). (In
terms of size).


This is how sizes of infinity are compared.

\begin{definition}
Suppose
\(T:\mathbb{R}^{n} \rightarrow \mathbb{R}^{m}\) is a linear map. What is
injective, surjective, and bijective? Injectivity is the same as
one-to-one. Bijective means isomorphism (one-to-one correspondence from
domain to codomain, for example, \(y = mx + b\) -- every single element
in the domain has a counterpart in the codomain, and the other way
around.)
\end{definition}


\section{\texorpdfstring{ Injective, onto, and
bijective}{ Injective, onto, and bijective}}

\begin{definition}
Suppose
\(T:\mathbb{R}^{n} \rightarrow \mathbb{R}^{m}\) is a linear map.
\end{definition}

\(T\) is \textbf{one to one} (injective) if:

\[T\left( \mathbf{v} \right) = T\left( \mathbf{u} \right) \Rightarrow \mathbf{u} = \mathbf{v}\]

Or if \(\mathbf{u} \neq \mathbf{v}\), then
\(T(\mathbf{u}) \neq T\left( \mathbf{v} \right)\).

For example: \(\left\{ \begin{matrix}
f(x) = x & \text{one\ to\ one} \\
f(x) = x^{2} & \text{one\ to\ many} \\
\end{matrix} \right.\ \)

\begin{quote}
\(\ker(T) = \left\{ \mathbf{0} \right\} \Rightarrow\) \(T\) is one
to one.
\end{quote}

A function is \textbf{onto} if:

\begin{enumerate}
\def\labelenumi{\roman{enumi})}
\item
  \(T\left( \mathbb{R}^{n} \right) = \mathbb{R}^{m}\) (The image/range
  covers the entire codomain)
\item
  \(\forall v' \in \mathbb{R}^{m},\ \exists v \in \mathbb{R}^{n}\ \text{s.t.}\text{\ T}(v) = v'\)
\item
  \(\text{Range}(T) = \mathbb{R}^{m}\)
\item
  \(T\left( \mathbf{x} \right) = A\mathbf{x}\), only if
  \(\text{col}(A) = \mathbb{R}^{m}\) and \(\text{rank}{(A)} = m\)
\end{enumerate}

These three statements are equivalent. (A function is onto if every
element on the codomain can be traced back to the domain)

Example: \(f(x) = x\) is onto, but \(f(x) = e^{x}\) is \textbf{not} an
onto function as outputs are strictly positive. If you try to reverse a
negative value (which is in the codomain) it does not exist.


\subsection{Isomorphism (Bijective):}

If a linear map is one-to-one and onto, the linear map is bijective.


\begin{remark}
Let \(T:\mathbb{R}^{n} \rightarrow \mathbb{R}^{m}\) be
a linear map. Represent \(T\) as
\(T\left( \mathbf{x} \right) = A\mathbf{x}\). The input vector
must be \(n\)-dimensional, so column count of \(A\) is \(n\). Output
vector must be \(m\)-dimensional, so row count of \(A\) is \(m\).
\end{remark}

To compute \(A\), it is the image of the basis of
\({\mathbf{e}}_{1}\ldots e_{n}\).

\[\begin{bmatrix}
 \mid & \mid & \mid & \cdots & \mid \\
T\left( {\mathbf{e}}_{1} \right) & T\left( {\mathbf{e}}_{2} \right) & T\left( {\mathbf{e}}_{3} \right) & \cdots & T\left( {\mathbf{e}}_{n} \right) \\
 \mid & \mid & \mid & \cdots & \mid \\
\end{bmatrix}\begin{bmatrix}
x_{1} \\
x_{2} \\
 \vdots \\
x_{n} \\
\end{bmatrix} = \begin{bmatrix}
y_{1} \\
 \vdots \\
 \vdots \\
y_{m} \\
\end{bmatrix}\]

\begin{framed}

\begin{enumerate}
\def\labelenumi{\arabic{enumi}.}
\item
  A function is one to one \(\Leftrightarrow\)
  \(\ker(T) = \text{null}(A) = \left\{ \mathbf{0} \right\}\) (The
  kernel of \(T\) is strictly the zero vector.) \textbf{THE COLUMNS OF}
  \(\mathbf{A}\) \textbf{ARE LINEARLY INDEPENDENT.}
\item
  A function is onto \(\Leftrightarrow\)
  \(\text{col}(A) = \mathbb{R}^{m}\) or \(\text{Rank}{(A)} = m\) (full
  rank.) (\textbf{COL SPACE MUST MATCH THE CODOMAIN or ALL ROWS HAVE
  PIVOTS or RANK MATCHES DIMENSION OF OUTPUT)}
\item
  Isomorphism: \(A\sim I \Leftrightarrow A^{- 1}\)
\end{enumerate}

\end{framed}



Example: \(T:\mathbb{R}^{3} \rightarrow \mathbb{R}^{3}\) is defined by
\(T\left( \begin{bmatrix}
x_{1} & x_{2} & x_{3} \\
\end{bmatrix} \right) = \left\lbrack x_{1} - 2x_{2},\ x_{2} + x_{3},\  - 2x_{3} - x_{1} \right\rbrack\)

Check if this is one-to-one, onto, or isomorphic.

For example:
\(T\left( \mathbf{x} + r\mathbf{y} \right) = T\left( \mathbf{x} \right) + rT\left( \mathbf{y} \right)\)

\begin{align*}
T\left( \mathbf{x} \right) &= A\mathbf{x} \\
T\left( \mathbf{x} + r\mathbf{y} \right) &= A\left( \mathbf{x} + r\mathbf{y} \right) = A\mathbf{x} + rA\mathbf{y}
\end{align*}


Example: \(H\mathbf{x} = 0\), \(\mathbf{x}\) unique?

\[\begin{bmatrix}
1 & 0 & 2 \\
0 & 1 & 1 \\
0 & 0 & 0 \\
\end{bmatrix}\begin{bmatrix}
x_{1} \\
x_{2} \\
x_{3} \\
\end{bmatrix} = \begin{bmatrix}
x_{1} + 2x_{3} \\
x_{2} + x_{3} \\
0 \\
\end{bmatrix}\]

This implies that \(x_{3}\) is a free variable, meaning anything in this
vector in the map brings it to zero:

\[s\begin{bmatrix}
 - 2 \\
 - 1 \\
1 \\
\end{bmatrix},\ s\mathbb{\in R}\]

This implies: \({\mathbf{u}}_{1} \neq {\mathbf{u}}_{2}\). i.e.,
\({\mathbf{u}}_{1} = \lbrack - 2,\  - 1,\ 1\rbrack,\ {\mathbf{u}}_{2} = \lbrack - 4,\ 2,\ 2\rbrack,\ T\left( {\mathbf{u}}_{1} \right) = 0\ \text{and}\text{\ T}\left( {\mathbf{u}}_{2} \right) = 0\)

Based on the kernel argument, you can find two different vectors,
meaning at least two different vectors map to zero.

Onto:
\(\dim\left( \text{codomain} \right) = \dim\left( \mathbb{R}^{3} \right) = 3\)

\[\text{rank}(A) = 2\]

(The rank is 2 because \(H\) only has two pivots.)

Isomorphic: NOT. \(A\) is not invertible.

\textbf{CHECK FOR ISOMORPHISM}

\[T:\mathbb{R}^{N} \rightarrow \mathbb{R}^{n},\ A \in M_{n \times n}\left( \mathbb{R}^{2} \right)\]

(The corresponding matrix has to be square)

\[A\mathbf{x} = \mathbf{0} \rightarrow \text{unique}\ \mathbf{x} = 0\]

(There should be a pivot for every column)

\[A\sim I_{n \times n}\]


\subsubsection{Example 1: One-to-one but not
onto}

\[T:\mathbb{R}^{3} \rightarrow \mathbb{R}^{4}\]

\[{T\left( \left\lbrack x_{1},\ x_{2},\ x_{3} \right\rbrack \right) }{= \left\lbrack x_{1} - x_{2} + x_{3},\ x_{2} - x_{3},\ x_{3} - x_{1},\ x_{1} + 2x_{2} + x_{3} \right\rbrack} \]

Check if one to one / onto / isomorphic.

Firstly, this transformation is not isomorphic as the input and output
dimension does not match. The size of \(A\) will be \(4 \times 3\), as
you expect the input to be a 3D vector meaning the matrix must contain
three columns.

The matrix \(A\): \(\dim\left( \text{col}(A) \right) = 3\), which is the
best case. The column space cannot cover \(\mathbb{R}^{4}\).

\(\dim\left( \text{col}(A) \right) = 3 \Rightarrow \text{Range}(T)\) is
strictly a subset of \(R^{4}\) and is NOT \(R^{4}\). The linear
transformation is not onto.


The map \textbf{can} be one-to-one: Just set up
\(A\mathbf{x} = \mathbf{0}\) and check the solution count. All
you need to do is check if it has free variables.

Write down the image of the standard basis:

\[A = \begin{bmatrix}
 \mid & \mid & \mid \\
T\left( {\mathbf{e}}_{1} \right) & T\left( {\mathbf{e}}_{2} \right) & T\left( {\mathbf{e}}_{3} \right) \\
 \mid & \mid & \mid \\
\end{bmatrix} = \begin{bmatrix}
1 & - 1 & 1 \\
0 & 1 & - 1 \\
 - 1 & 0 & 1 \\
1 & 2 & 1 \\
\end{bmatrix}\sim\begin{bmatrix}
1 & 0 & 0 \\
0 & 1 & 0 \\
0 & 0 & 1 \\
0 & 0 & 0 \\
\end{bmatrix}\]

No free variable. We have a unique solution for
\(A\mathbf{x} = \mathbf{0}\), meaning \(T\) is one-to-one as
\(\ker(T) = null(A) = \left\{ \mathbf{0} \right\}\).


\subsubsection{Example 2: Onto but not
one-to-one}

\[{T:\mathbb{R}^{3} \rightarrow \mathbb{R}^{2};T\left( \left\lbrack x_{1},\ x_{2},\ x_{3} \right\rbrack \right) }{= \left\lbrack x_{1} - 2x_{2} + 3x_{3},\ x_{2} - 2x_{3} \right\rbrack} \]

Matrix \(A \in M_{2 \times 3}\)


The matrix being wider than it is tall → \(\exists\) free variable. The
system \(A\mathbf{x} = \mathbf{0}\) cannot give you a unique
solution, implying infinitely as many solutions. This means

\[\ker(T) = \text{Null}(A) \neq \left\{ \mathbf{0} \right\} \Rightarrow \text{T\ is\ not\ one-to-one}\]

Check for onto:
\(\text{Rank}(A) = 2 \Rightarrow \text{col}(A) = \mathbb{R}^{2}\)

\[A = \begin{bmatrix}
1 & - 2 & 3 \\
0 & 1 & - 2 \\
\end{bmatrix}\sim\begin{bmatrix}
1 & 0 & - 1 \\
0 & 1 & - 2 \\
\end{bmatrix}\]

We know every single \textbf{row} contains a pivot.


\subsubsection{Example 3:}

\[T:\mathbb{R}^{3} \rightarrow \mathbb{R}^{3},\ A \in M_{3 \times 3}\]

\[{T:\left( \left\lbrack x_{1},\ x_{2},\ x_{3} \right\rbrack \right) }{= \left\lbrack x_{2} - 2x_{2} + 3x_{3},\ x_{2} - 2x_{3},\ x_{1} + x_{2} + x_{3} \right\rbrack} \]

If the domain and the codomain have the same dimension computations have
to be done to figure out if the transformation is one-to-one, onto, or
isomorphic.

Solution: \(A = \begin{bmatrix}
1 & - 2 & 3 \\
0 & 1 & - 2 \\
1 & 1 & 1 \\
\end{bmatrix}\sim\begin{bmatrix}
1 & - 2 & - 3 \\
0 & 1 & - 2 \\
0 & 0 & 4 \\
\end{bmatrix}\sim I\)

This matrix is not RREF but can convince me that the matrix can be
reduced to an identity matrix.

Therefore, the linear system is \textbf{isomorphic}, implying it is
one-to-one and onto. (\(T\) is isomorphic).

\begin{align*}
A\sim I \Rightarrow Rank(A) &= 3 = \dim\left( \mathbb{R}^{3} \right) \Rightarrow \text{onto} \\
&\text{No\ free\ variable} \Rightarrow A\mathbf{x}  \\
&= \mathbf{0}\ \text{has\ a\ unique\ solution} \Rightarrow null(A)  \\
&= \left\{ \mathbf{0} \right) \Leftrightarrow \ker(T) = \left\{ \mathbf{0} \right\}  \\
&\Rightarrow T\ \text{is\ one-to-one} \Rightarrow T\ \text{is\ isomorphic.}
\end{align*}


(Longer explanation above)

\begin{framed}

\textbf{One-to-one but not onto} \(\mathbf{\Rightarrow}\) \textbf{The
matrix is taller than it is wide.} \(\begin{bmatrix}
a & b \\
c & d \\
e & f \\
\end{bmatrix}\)

\begin{itemize}
\item
  \textbf{Typically} has pivots in every column \(\Leftrightarrow\)
  column space is typically linearly independent.
\item
  \textbf{NO FREE VARIABLES}
\item
  Will not have pivots in every row.
\item
  Column space does not span the codomain. For example, in a
  \(3 \times 2\) matrix, \(\dim\left( \text{col}(A) \right) \leq 2\). It
  needs to be 3 to be onto.
\item
  If you have a 2D input your image will only be a plane, or something
  that is worse than that, unless you try to compress a 3D space into
  2D, in which you can't in this case.
\end{itemize}

\textbf{Onto but not one-to-one} \(\mathbf{\Rightarrow}\) \textbf{The
matrix is wider than it is tall.} \(\begin{bmatrix}
a & b & c \\
d & e & f \\
\end{bmatrix}\)

\begin{itemize}
\item
  In a \(2 \times 3\) matrix, \(\text{rank}(A) \leq 2\). If
  \(\text{rank}(A) = 2\), which occurs most of the time, the matrix is
  onto. \textbf{(COLUMN PIVOT COUNT = HEIGHT OF MATRIX)}
\item
  \(\text{rank}(A) \leq 2 \Rightarrow \text{rank}(A) < 3\). The rank can
  never match the width of the matrix \(\Leftrightarrow\) has pivotless
  columns \(\Leftrightarrow\) matrix has free variable
  \(\Leftrightarrow\) the span of the column space has redundant
  vectors.
\item
  Compressing a 3D input to a 2D output, of course the existence of one
  of the entries in the 4D input will be redundant.
\end{itemize}

\textbf{Neither onto nor one-to-one:}

\begin{itemize}
\item
  Read the above two like a flow chart. There are conditions on the
  first point. If the condition fails, it is neither.
\end{itemize}

\textbf{Square matrices:}

For square matrix \(A \in M_{n \times n}\), \(A\) is one-to-one
\(\Leftrightarrow\) \(A\) is onto. \textbf{(A SQUARE MATRIX IS EITHER
ISOMORPHIC, OR NEITHER ONE-TO-ONE OR ONTO.)}

RATIONALE: one-to-one requires pivots in all columns; onto requires
pivots in every row. These notions become equivalent if the matrix is
square.

\textbf{Isomorphism:}

\(A\sim I \Leftrightarrow A\) is invertible \(\Leftrightarrow\) \(A\) is
one-to-one and \(A\) is onto.

\end{framed}




\section{Determinant}

Represented by \(\det(A)\) or \(|A|\) (this isn't absolute value)

In lower dimensional cases
(\(\mathbb{R}^{1},\ \mathbb{R}^{2},\ \mathbb{R}^{3})\):

\begin{definition}
The determinant of a \(1 \times 1\) matrix is just
its entry.

\[A \in M_{1 \times 1} \Rightarrow \det(A) = \alpha\mathbb{\in R}\]
\end{definition}

The determinant of a \(2 \times 2\) matrix is:

\[A = \begin{bmatrix}
a_{1}\  & a_{2} \\
b_{1} & b_{2} \\
\end{bmatrix} \Rightarrow \det(A) = \left| \begin{bmatrix}
a_{1}\  & a_{2} \\
b_{1} & b_{2} \\
\end{bmatrix} \right| = a_{1}b_{2} - a_{2}b_{1}\]

The determinant of a \(3 \times 3\) matrix is:

\[A = \begin{bmatrix}
a_{1} & a_{2} & a_{3} \\
b_{1} & b_{2} & b_{3} \\
c_{1} & c_{2} & c_{3} \\
\end{bmatrix}\]

\[{\det(A) = \left| \begin{matrix}
a_{1} & a_{2} & a_{3} \\
b_{1} & b_{2} & b_{3} \\
c_{1} & c_{2} & c_{3} \\
\end{matrix} \right| }{= a_{1}\left| \begin{matrix}
b_{2} & b_{3} \\
c_{2} & c_{3} \\
\end{matrix} \right| - a_{2}\left| \begin{matrix}
b_{1} & b_{3} \\
c_{1} & c_{3} \\
\end{matrix} \right| }{+ a_{3}\left| \begin{matrix}
b_{1} & b_{2} \\
c_{1} & c_{2} \\
\end{matrix} \right|} \]


\subsection{Cross product}

If
\(\mathbf{a} = \left\lbrack a_{1},\ a_{2},\ a_{3} \right\rbrack \in \mathbb{R}^{3},\ \mathbf{b} = \left\lbrack b_{1},\ b_{2},\ b_{3} \right\rbrack \in \mathbb{R}^{3}\)

\[\mathbf{a} \times \mathbf{b} = \det\left| \begin{matrix}
i & j & k \\
a_{1} & a_{2} & a_{3} \\
b_{1} & b_{2} & b_{3} \\
\end{matrix} \right| = i\left| \begin{matrix}
a_{2} & a_{3} \\
b_{2} & b_{3} \\
\end{matrix} \right| - j\left| \begin{matrix}
a_{1} & a_{3} \\
b_{1} & b_{3} \\
\end{matrix} \right| + k\left| \begin{matrix}
a_{1} & a_{2} \\
b_{1} & b_{2} \\
\end{matrix} \right|\]

In vector form:

\[\mathbf{v} = \left( \left| \begin{matrix}
a_{2} & a_{3} \\
b_{2} & b_{3} \\
\end{matrix} \right|,\  - \left| \begin{matrix}
a_{1} & a_{3} \\
b_{1} & b_{3} \\
\end{matrix} \right|,\ \left| \begin{matrix}
a_{1} & a_{2} \\
b_{1} & b_{2} \\
\end{matrix} \right| \right)\]

We also want to check if
\(\mathbf{v}\bot\mathbf{a},\ \mathbf{v}\bot\mathbf{b}\).
Just use the dot product.

\begin{align*}
\mathbf{a} \cdot \mathbf{v} &= \left( a_{1},\ a_{2},\ a_{3} \right) \cdot \left( \mathbf{a} \times \mathbf{b} \right) \\
&= \left( a_{1},\ a_{2},\ a_{3} \right) \cdot \left( \left| \begin{matrix}a_{2} & a_{3} \\b_{2} & b_{3} \\\end{matrix} \right|,\  - \left| \begin{matrix}a_{1} & a_{3} \\b_{1} & b_{3} \\\end{matrix} \right|,\ \left| \begin{matrix}a_{1} & a_{2} \\b_{1} & b_{2} \\\end{matrix} \right| \right) \\
& \\
&= \left( a_{1},\ a_{2},\ a_{3} \right) \cdot \left( a_{2}b_{3} - a_{3}b_{2},\  - a_{1}b_{3} + a_{3}b_{1},\ a_{1}b_{2} - a_{2}b_{1} \right) \\
& \\
&= a_{1}a_{2}b_{3} - a_{1}a_{3}b_{2} - a_{2}a_{1}b_{3} + a_{2}a_{3}b_{1}  \\
&+ a_{3}a_{1}b_{2} - a_{3}a_{2}b_{1} \\
&= 0
\end{align*}


(All cancels out. Similar computations can be done for \(b\).)

\textbf{Lemma:}
\(\mathbf{a} \times \mathbf{b} = \mathbf{v}\) such that
\(\mathbf{v}\bot\mathbf{a},\ \mathbf{v}\bot\mathbf{b}\)


\subsection{Geometric interpretation of
determinants}

Determinant of a \(2 \times 2\) matrix

\[\begin{bmatrix}
 - & {\mathbf{v}}_{1} & - \\
 - & {\mathbf{v}}_{2} & - \\
\end{bmatrix},\ {\mathbf{v}}_{1},\ {\mathbf{v}}_{2} \in \mathbb{R}^{2}\]

\[\det(A) = \text{The\ area\ defined\ by\ }{\mathbf{v}}_{1},\ {\mathbf{v}}_{2}\]


If \({\mathbf{v}}_{1} \parallel {\mathbf{v}}_{2}\), the area of
the parallelogram formed by it is 0, so \(\det(A) = 0\).

\textbf{Proof.}

\[{\text{Area}^{2} = \left( \left| \left| \mathbf{a} \right| \right|\left| \left| \mathbf{h} \right| \right| \right)^{2} }{= \left( \left| \left| \mathbf{a} \right| \right| \cdot \left| \left| \mathbf{b} \right| \right| \cdot \sin^{2}\theta \right) }{= \left( \left| \left| \mathbf{a} \right| \right| \cdot \left| \left| \mathbf{b} \right| \right| \cdot (1 - \cos^{2}\theta \right) }{= \cdots} \]


Dealing with three-dimensional vectors:

\[\mathbf{a},\ \mathbf{b},\ \mathbf{c} \in \mathbb{R}^{3}\]

If the three vectors are mutually independent, a parallelepiped can be
formed.


The determinant of the three vectors tells us the volume of this thing.

The case where \(\mathbf{a},\ \mathbf{b},\ \mathbf{c}\) are
linearly dependent, the volume of the parallelepiped is zero.


If all vectors are parallel to each other, they define a 1-dimensional
object.


Either case, \(\det(A) = 0,\) where \(A = \begin{bmatrix}
a_{1} & a_{2} & a_{3} \\
b_{1} & b_{2} & b_{3} \\
c_{1} & c_{2} & c_{3} \\
\end{bmatrix}\)

\textbf{Example:} All these operations have the same intent.

Find the volume of the parallelepiped determined by
\(\mathbf{a} = \lbrack 1,\ 0,\ 3\rbrack,\ \mathbf{b} = \lbrack 0,\ 1,\ 2\rbrack,\ \mathbf{c} = \lbrack 3,\ 3,\ 1\rbrack\)
and check if it isn't 0

\(\Leftrightarrow \ \)Check of
\(\mathbf{a},\ \mathbf{b},\mathbf{c}\) are linearly
independent

\(\Leftrightarrow\) Check if
\(\left\{ \mathbf{a},\ \mathbf{b},\ \mathbf{c} \right\}\) is
a basis for \(\mathbb{R}^{3}\)

\(\Leftrightarrow\) Check if \(A\mathbf{x} = \mathbf{d}\) has a
unique solution where \(A = \begin{bmatrix}
a_{1} & a_{2} & a_{3} \\
b_{1} & b_{2} & b_{3} \\
c_{1} & c_{2} & c_{3} \\
\end{bmatrix}\)

\(\Leftrightarrow\) Check of \(T = A\mathbf{x}\) is isomorphic,
one-to-one, or onto

\(\Leftrightarrow\) check if \(A\) is invertible.

\textbf{Solution:} Compute the determinant.

\[\text{Vol} = \det(A) = \left| \begin{matrix}
1 & 0 & 3 \\
0 & 1 & 2 \\
3 & 3 & 1 \\
\end{matrix} \right| = 1 \cdot \left| \begin{matrix}
1 & 2 \\
3 & 1 \\
\end{matrix} \right| - 0 + 3 \cdot \left| \begin{matrix}
0 & 1 \\
3 & 3 \\
\end{matrix} \right| = | - 14| = 14\]

The volume must be positive.

The position of the vectors must form something like this.


Therefore, these vectors are linearly independent, so they can span 3D
space. The matrix \(A\sim I\), has no free variable, the equation
\(A\mathbf{x} = \mathbf{d}\) has a unique solution, is
full-rank, and the matrix is invertible.

\textbf{THE VOLUME OF A PARALLELPIPED FROM THREE ADJACENT VECTORS IS THE
DETERMINANT OF THE MATRIX FORMED BY THE ROW VECTORS}


\subsection{Generalization of
determinant}

The row vectors determine the edge of the polyhedron formed by the
vectors. If the polyhedron has a nonzero volume, the row vectors are
linearly independent.

\textbf{Definition:} The minor matrix \(A_{ij}\) for
\(A \in M_{n \times n}\left( \mathbb{R} \right)\) is the
\((n - 1) \times (n - 1)\) matrix. Obtain this matrix by removing the
\(i\)th row and \(j\)th column of \(A\) (remove the column and row for
the element we are targeting).

Example: The minor of \(b_{3}\) is

\[\begin{bmatrix}
a_{1} & a_{2} & a_{3} & a_{4} \\
b_{1} & b_{2} & b_{3} & b_{4} \\
c_{1} & c_{2} & c_{3} & c_{4} \\
d_{1} & d_{2} & d_{3} & d_{4} \\
\end{bmatrix} \rightarrow \begin{bmatrix}
a_{1} & a_{2} & a_{4} \\
c_{1} & c_{2} & c_{4} \\
d_{1} & d_{2} & d_{4} \\
\end{bmatrix}\]

The cofactor of \(a_{ij}\) of \(A\) is:

\[{a_{ij}' }{= ( - 1)^{1 + j}\det\left( A_{ij} \right)\ \text{where\ }A_{ij}\ \text{is\ the\ minor\ for\ }a_{ij}} \]

Based on these two, we are able to determine the determinant for an
\(n \times n\) matrix \(A\).

\[{\det(A) = \left| \begin{matrix}
a_{11} & a_{12} & \cdots & a_{1n} \\
a_{21} & \ddots & \ddots & a_{2n} \\
 \vdots & \ddots & \ddots & \vdots \\
a_{n1} & a_{n2} & \cdots & a_{nn} \\
\end{matrix} \right| }{= a_{11} \cdot a_{11}' + a_{12} \cdot a_{12}' + \cdots + a_{1n} \cdot a_{1n}'} \]

For example:

\[{\text{d}et}\begin{bmatrix}
a_{1} & a_{2} & a_{3} & a_{4} \\
b_{1} & b_{2} & b_{3} & b_{4} \\
c_{1} & c_{2} & c_{3} & c_{4} \\
d_{1} & d_{2} & d_{3} & d_{4} \\
\end{bmatrix}\]

\[= a_{1}a_{1}' + a_{2}a_{2}' + a_{3}a_{3}' + a_{4}a_{4}'\]

\[{= a_{1}( - 1)^{1 + 1}\det\left( \begin{bmatrix}
b_{1} & b_{3} & b_{4} \\
c_{2} & c_{3} & c_{4} \\
d_{2} & d_{3} & d_{4} \\
\end{bmatrix} \right) }{+ a_{2}( - 1)^{1 + 2}\det\left( \begin{bmatrix}
b_{1} & b_{3} & b_{4} \\
c_{1} & c_{3} & c_{4} \\
d_{1} & d_{3} & d_{4} \\
\end{bmatrix} \right) }{+ a_{3}( - 1)^{1 + 3}\det\left( \begin{bmatrix}
b_{1} & b_{2} & b_{4} \\
c_{1} & c_{2} & c_{4} \\
d_{1} & d_{2} & d_{4} \\
\end{bmatrix} \right) }{+ a_{4}( - 1)^{1 + 4}\det\left( \begin{bmatrix}
b_{1} & b_{2} & b_{3} \\
c_{1} & c_{2} & c_{3} \\
d_{1} & d_{2} & d_{3} \\
\end{bmatrix} \right)} \]

The determinant of a \(5 \times 5\) matrix:

\begin{align*}
\det(A) &= a_{1}a_{1}' + a_{2}a_{2}' + \cdots + a_{5}a_{5}' \\
&= a_{1}( - 1)^{1 + 1}\det(4 \times 4) + \cdots + a_{5}( - 1)^{1 + 5}\det(4 \times 4)
\end{align*}


To compute the determinant for a \(5 \times 5\) matrix, you are required
to compute the determinant of a \(4 \times 4\) matrix 5 times.

\textbf{You can work with any row you want as long as the rows stay the
same.}

Computing these are impossible by hand.

(You do not need to prove properties of determinants / use the
properties to make computations easier)

\(\text{adj}(A) = C^{T}\), where \(C\) is the minor matrix of \(A\),
a.k.a. \(C = \begin{bmatrix}
a_{11}' & a_{12}' \\
a_{21}' & a_{22}' \\
\end{bmatrix}\) if \(A \in M_{2 \times 2}\)


\subsection{Tricks for computing
determinant}

\begin{itemize}
\item
  Compute \(\det(A)\) along the row with the most zeroes
\item
  Compute \(\det\left( A^{T} \right)\) if a column has many zeroes
  (rotate)
\end{itemize}


\subsubsection{Targeting different
rows}

\emph{Our target row is the row we are working with.}

When calculating the determinant of a larger matrix, aim to compute it
along the row with the most zeroes. For example:

\[A = \begin{bmatrix}
1 & 3 & - 2 \\
0 & 1 & 5 \\
 - 2 & - 6 & 7 \\
\end{bmatrix},\det(A) = 1 \cdot \left| \begin{matrix}
1 & 5 \\
 - 6 & 7 \\
\end{matrix} \right| - 3\left| \begin{matrix}
0 & 5 \\
 - 2 & 7 \\
\end{matrix} \right| - 2\left| \begin{matrix}
0 & 1 \\
 - 2 & 6 \\
\end{matrix} \right|\]

You'll have to compute \(2 \times 2\) determinants three times. However,
when we target the second row:

\[\det(A) = 0( - 1)^{2 + 1}\left| \begin{matrix}
3 & - 2 \\
 - 6 & 7 \\
\end{matrix} \right| + 1( - 1)^{2 + 2}\left| \begin{matrix}
1 & - 2 \\
 - 2 & 7 \\
\end{matrix} \right| + 5( - 1)^{2 + 3}\left| \begin{matrix}
1 & 3 \\
2 & - 6 \\
\end{matrix} \right|\]

You will notice that one of the terms added is multiplied by zero, so we
can get rid of that:

\[\det(A) = 1( - 1)^{2 + 2}\left| \begin{matrix}
1 & - 2 \\
 - 2 & 7 \\
\end{matrix} \right| + 5( - 1)^{2 + 3}\left| \begin{matrix}
1 & 3 \\
2 & - 6 \\
\end{matrix} \right|\]


\subsubsection{Transpose invariance}

\[\det(A) = \det\left( A^{T} \right)\]

Working with the columns is the same as working with a row.

Example: \(A = \begin{bmatrix}
1 & 2 & 0 & 0 \\
0 & 4 & 1 & 5 \\
0 & 2 & 2 & 6 \\
0 & 1 & 4 & 7 \\
\end{bmatrix},\det(A) = \det\left( A^{T} \right) = \left| \begin{matrix}
1 & 0 & 0 & 0 \\
2 & 4 & 2 & 1 \\
0 & 1 & 2 & 4 \\
0 & 5 & 6 & 7 \\
\end{matrix} \right| = 1( - 1)^{1 + 1}\left| \begin{matrix}
4 & 2 & 1 \\
1 & 2 & 4 \\
5 & 6 & 7 \\
\end{matrix} \right|\)


\subsubsection{Triangular and diagonal
matrices}

If \(A\) is a triangular matrix (either upper or lower triangular), then
\(\det(A) =\) product of all its diagonal entries.

Example:

\[{A = \begin{bmatrix}
a_{11} & a_{12} & a_{12} \\
0 & a_{22} & a_{23} \\
0 & 0 & a_{33} \\
\end{bmatrix},\det(A) }{= a_{33} \cdot a_{33}' }{= a_{33}( - 1)^{3 + 3}\left| \begin{matrix}
a_{11} & a_{12} \\
0 & a_{22} \\
\end{matrix} \right| }{= a_{11} \cdot a_{22} \cdot a_{33}} \]

\[{\left| \begin{matrix}
a_{11} & * & * & * \\
0 & a_{22} & * & * \\
0 & 0 & \ddots & * \\
0 & 0 & 0 & a_{nn} \\
\end{matrix} \right| }{= a_{nn}( - 1)^{n + n}\left| \begin{matrix}
a_{11} & * & * \\
0 & \ddots & * \\
0 & 0 & a_{(n - 1)(n - 1)} \\
\end{matrix} \right| }{= \cdots = a_{nn}a_{(n - 1)(n - 1)}\cdots a_{11}} \]

If one of the diagonal entries are zero for an upper triangular matrix,
then the determinant is zero.

The same applies to diagonal matrices. \(\left| \begin{matrix}
2 & 0 & 0 \\
0 & 3 & 0 \\
0 & 0 & 4 \\
\end{matrix} \right| = 2 \cdot 3 \cdot 4\)


\subsubsection{Row operations}

\begin{enumerate}
\def\labelenumi{\arabic{enumi}.}
\item
  \textbf{ROW SWAPS NEGATE:} Let \(A \in M_{n \times n}\) matrix.
  Perform the row operation \(R_{i} \leftrightarrow R_{j}\). Let \(A'\)
  be the \(n \times n\) matrix after row interchange of one time. Then
  \(\det(A) = - \det\left( A' \right)\).
\item
  \textbf{LINEAR DEPENDENCE OF COLUMNS/ROWS MAKES IT ZERO:} If \(A\) has
  two equal or linearly dependent rows, then \(\det(A) = 0\) (implies
  the column space of the matrix is linearly dependent and compresses
  everything to a lower dimension, which always makes the volume zero)
\item
  \textbf{ROW SCALING ALSO SCALES THE DETERMINANT:}
  \(R_{i} \rightarrow r \cdot R_{i}\), \(\det(A) = r \cdot \det(A)\)
  \emph{Scaling one row of a matrix scales the determinant of the matrix
  by how much you scaled that row.}
\item
  \textbf{SCALING THE ENTIRE MATRIX APPLIES THE SQUARE-CUBE LAW:}
  \(r \cdot A = A' \Rightarrow \det\left( A' \right) = r^{n}\det(A)\)
  \emph{(Based off the square-cube law -- if a} \(2 \times 2\)
  \emph{matrix scales everything by two the determinant is 4 or the
  square of 2)}
\item
  \textbf{DETERMINANT WON'T CHANGE FOR ROW ADDITIONS:}
  \(R_{i} \rightarrow R_{i} + rR_{j}\),
  \(\det(A) = \det\left( A' \right)\)
\end{enumerate}

\emph{Note: The determinant of a matrix's RREF is not necessarily the
same as the original matrix, as row scaling could have been done.}

\[A = \begin{bmatrix}
1 & 3 & - 2 \\
0 & 1 & 5 \\
 - 2 & - 6 & 7 \\
\end{bmatrix}\sim A' = \begin{bmatrix}
1 & 3 & - 2 \\
0 & 1 & 5 \\
 - 2 & - 5 & 12 \\
\end{bmatrix}\]

\[R_{3} + R_{2};\det(A) = \det\left( A' \right) = 3\]

\[R_{3} - R_{2},\det\left( A'' \right) = 3\]

\begin{enumerate}
\def\labelenumi{\arabic{enumi}.}
\setcounter{enumi}{5}
\item
  \textbf{SPLITTING MATRIX MULTIPLICATION:} Let
  \(A,\ B \in M_{n \times n}\).
  \(\det\left( AB \right) = \det(A) \cdot \det(B)\)
\item
  \textbf{DETERMINANT OF} \(\mathbf{A}^{\mathbf{- 1}}\)\textbf{:} If
  \(A\) is invertible, then
  \(\det\left( A^{- 1} \right) = \frac{1}{\det(A)}\)
\end{enumerate}

Remark: \(A\) is invertible
\(\Rightarrow \exists A^{- 1},\ A \cdot A^{- 1} = A^{- 1} \cdot A = I\)

\begin{align*}
\det(I) &= \det\left( A \cdot A^{- 1} \right) = \det(A) \cdot \det\left( A^{- 1} \right) = 1 \\
1 &= \det(A) \cdot \det\left( A^{- 1} \right) \Rightarrow \frac{1}{\det(A)} = \det\left( A^{- 1} \right)
\end{align*}



\section{Eigenvalues, Eigenvectors}

\begin{definition}
Let
\(A \in M_{n \times n}\left( \mathbb{R} \right)\) (square matrix). A
scalar \(\lambda\) is an eigenvalue of \(A\) if there exists a non-zero
vector \(\mathbf{v}\backslash n\ \mathbb{R}^{n}\) such that
\(A\mathbf{v} = \lambda\mathbf{v}\). We will call
\(\mathbf{v}\) the eigenvector of \(A\) corresponding to the
eigenvalue \(\lambda\) (eigenvectors must correspond to a specific
eigenvalue).
\end{definition}


\subsection{Computation of
eigenvalues}

To find the eigenvalues of \(A\), I am looking for \(\lambda\) such that
\(\lambda\mathbf{v} = A\mathbf{v}\), where \(\mathbf{v}\) is
a corresponding eigenvector, where \(\mathbf{v} \neq 0\). We can
call the pair \(\left( \lambda,\ \mathbf{v} \right)\).

\begin{align*}
\lambda\mathbf{v} &= A\mathbf{v} \\
\Leftrightarrow 0 &= A\mathbf{v} - \lambda\mathbf{v} \\
\Leftrightarrow 0 &= (A - I\lambda)\mathbf{v} \\
\Leftrightarrow \det(A - I\lambda) &= 0
\end{align*}


\(\mathbf{v} \neq 0\) is a solution for
\(\left( A - I\lambda  \right)\). However, because
\(\left( A - I\lambda  \right)\mathbf{v} = \mathbf{0}\), then
\(\left( A - I\lambda  \right)\)'s determinant is zero.

\(\Rightarrow \lambda\) is an eigenvalue of \(A\) if it sastifies
\(\det(A - \lambda I) = 0\).

Recall: \(A\mathbf{x} = \mathbf{0} \Rightarrow\) unique solution
where \(\mathbf{x} = \mathbf{0}\).


\subsection{Computation of eigenspaces and
eigenvectors}

\begin{definition}
Let
\(A \in M_{n \times n}\left( \mathbb{R} \right)\). The characteristic
polynomial of \(A\) is \(p(\lambda) = \det(A - \lambda I)\) (this
expression itself, either side of the equality is the polynomial). The
roots of this characteristic polynomial gives you the eigenvalues. The
set
\(E_{\lambda} = \left\{ \mathbf{x} \in \mathbb{R}^{n}\mid A\mathbf{x} = \lambda\mathbf{x} \right\}\)
is the eigenspace of \(\lambda\). This space contains the zero vector
and all the eigenvectors of \(A\) corresponding to \(\lambda\).

\[E_{\lambda} = \text{null}\left( A - \lambda I \right)\]
\end{definition}

Why?

\begin{align*}
E_{\lambda} &= \left\{ \mathbf{x} \in \mathbb{R}^{n}\mid A\mathbf{x} = \lambda\mathbf{x} \right\} \\
\Leftrightarrow E_{\lambda} &= \left\{ \mathbf{x} \in \mathbb{R}^{n}\mid A\mathbf{x} - \lambda\mathbf{x} = 0 \right\} \\
&= \left\{ \mathbf{x} \in \mathbb{R}^{n}\mid\left( A - \lambda I \right)\mathbf{x} = 0 \right\}
\end{align*}



\subsubsection{Computational example}

Suppose I have the square matrix \(A = \begin{bmatrix}
2 & 0 & 0 \\
1 & 2 & - 1 \\
1 & 3 & 2 \\
\end{bmatrix}\) and I want to compute the eigenvalues and the
eigenspace.

\begin{enumerate}
\def\labelenumi{\arabic{enumi}.}
\item
  Compute all eigenvalues for \(A\).
\item
  For each eigenvalue \(\lambda\) of \(A\), find its eigenspace
  \(E_{\lambda}\).
\end{enumerate}

\begin{framed}

\textbf{FINDING EIGENVALUES}

Find all \(\lambda\) for \(0 = \det(A - I\lambda)\).

\textbf{FINDING EIGENSPACES}

For each eigenvalue, compute the nullspace of \(A - I\lambda\).

The eigenvectors (nonzero) corresponding to that \(\lambda\) is spanned
by the nullspace of \(A - I\lambda\).

\end{framed}




\paragraph{Determining eigenvalues}

Solution: \(p(\lambda) = 0 \rightarrow\) roots of \(p(\lambda)\) gives
you all eigenvalues of \(A\).

\[0 = p(\lambda) = \det(A - \lambda I) = \left| \begin{matrix}
2 - \lambda & 0 & 0 \\
1 & 2 - \lambda & - 1 \\
1 & 3 & - 2 - \lambda \\
\end{matrix} \right|\]

\begin{align*}
&= (2 - \lambda)( - 1)^{1 + 1}\left| \begin{matrix}2 - \lambda & - 1 \\3 & - 2 - \lambda \\\end{matrix} \right| \\
&= (2 - \lambda)\left( (2 - \lambda)( - 2 - \lambda) + 3 \right) \\
&= (2 - \lambda)\left( \lambda^{2} - 1 \right) = (2 - \lambda)(\lambda - 1)(\lambda + 1)
\end{align*}


Based on this factorization
\(p(\lambda) = 0 \Rightarrow \lambda = 2,\ \lambda = 1,\ \lambda = - 1\)

Simplifying this gives you a polynomial to solve, where it is equated to
zero.

The eigenvalues of \(A\) gives 2, 1, and \(- 1\). We have three
different eigenspaces:

\[E_{\lambda = 2},\ E_{\lambda = 1},\ E_{\lambda = - 1}\]


\paragraph{Determining eigenspaces}

For \(\lambda = 2\), compute the nullspace of
\(A - \lambda I = A - 2I = \begin{bmatrix}
0 & 0 & 0 \\
1 & 0 & - 1 \\
1 & 3 & - 4 \\
\end{bmatrix}\). Nullspace:
\((A - \lambda I)\mathbf{x} = \mathbf{0}\)

Reduce the matrix. We will obtain \(\begin{bmatrix}
1 & 0 & - 1 \\
0 & 1 & - 1 \\
0 & 0 & 0 \\
\end{bmatrix}\). (We have \(x_{2} - x_{3} = 0\) and
\(x_{1} - x_{3} = 0\))

This reduced matrix has precisely one free variable, so
\(x_{3} = s,\ x_{2} = s,\ x_{1} = s\), so the general solution is
\(\mathbf{x} = \begin{bmatrix}
s \\
s \\
s \\
\end{bmatrix},\ s\mathbb{\in R}\). Therefore, the set of vectors in
\(E_{\lambda = 2} = \left\{ \mathbf{x} \in s\begin{bmatrix}
1 \\
1 \\
1 \\
\end{bmatrix},\ s\mathbb{\in R} \right\}\) (A set of vectors spanned by
\(\begin{bmatrix}
1 \\
1 \\
1 \\
\end{bmatrix}\) and has a subspace structure, as it is generated by span
meaning closed under addition and scalar multiplication.)

\[{\mathbf{b}}_{1} = \begin{bmatrix}
1 \\
1 \\
1 \\
\end{bmatrix} \rightarrow \text{eigenv}\text{ector}\]

An eigenvector corresponding to \(\lambda = 2\) is where
\(A\mathbf{v} = \lambda\mathbf{v}\).

\[\begin{matrix}
A\mathbf{v} = \begin{bmatrix}
2 & 0 & 0 \\
1 & 2 & - 1 \\
1 & 3 & - 2 \\
\end{bmatrix} \cdot \begin{bmatrix}
1 \\
1 \\
1 \\
\end{bmatrix} = \begin{bmatrix}
2 \\
2 \\
2 \\
\end{bmatrix} \\
\lambda\mathbf{v} = 2 \cdot \begin{bmatrix}
2 \\
2 \\
2 \\
\end{bmatrix} \\
\end{matrix} \Rightarrow A\mathbf{v} = \lambda\mathbf{v}\]

For \(E_{\lambda = 1}\):
\(\text{null}(A - 1I) = null\left( \begin{bmatrix}
1 & 0 & 0 \\
1 & 1 & - 1 \\
1 & 3 & 3 \\
\end{bmatrix} \right)\)

The nullspace is \(\mathbf{x}\) such that
\(\beta\mathbf{x} = \mathbf{0}\)

\[\begin{bmatrix}
1 & 0 & 0 \\
1 & 1 & - 1 \\
1 & 3 & 3 \\
\end{bmatrix}\sim\begin{bmatrix}
1 & 0 & 0 \\
0 & 1 & - 1 \\
0 & 0 & 0 \\
\end{bmatrix}\]

Therefore \(x_{3} = s,\ x_{2} = s,\ x_{1} = 0\) and
\(\mathbf{x} = sp\left\{ \begin{bmatrix}
0 \\
1 \\
1 \\
\end{bmatrix} \right\}\)

This means
\(E_{\lambda = 1} = \left\{ \mathbf{x} = sp\begin{bmatrix}
0 \\
1 \\
1 \\
\end{bmatrix} \right\}\)

This means \(A\mathbf{v} = \mathbf{\lambda}\) given
\(\mathbf{v}\) is in the span of \(\begin{bmatrix}
0 \\
1 \\
1 \\
\end{bmatrix}\).

For \(\lambda = - 1\)

\[\text{null}(A + I) = null\begin{bmatrix}
3 & 0 & 0 \\
1 & 3 & - 1 \\
1 & 3 & - 1 \\
\end{bmatrix}\sim null\begin{bmatrix}
1 & 0 & 0 \\
0 & 3 & - 1 \\
0 & 0 & 0 \\
\end{bmatrix}\]

To compute the solution, reduce the matrix.

\[\mathbf{x} = sp\begin{bmatrix}
0 \\
1 \\
3 \\
\end{bmatrix}\]

This is one concrete example for computation of eigenvalues and
corresponding eigenspaces. The eigenvectors are nonzero vectors in the
corresponding eigenspaces.


\subsection{Properties of eigenvalues}

\begin{enumerate}
\def\labelenumi{\arabic{enumi}.}
\item
  \(A^{k}\mathbf{v} = \lambda^{k}\mathbf{v}\) (Proof:
  \(A^{2}\mathbf{v} = AA\mathbf{v} = A\lambda \mathbf{v} = \lambda A\mathbf{v} = \lambda \lambda \mathbf{v} = \lambda^{2}\mathbf{v}\),
  then proceed with induction)
\item
  Zero (0) cannot be an eigenvalue for an invertible matrix.
\end{enumerate}

\begin{proof} \(A\) is invertible
\(\Leftrightarrow \det(A) \neq 0 \Leftrightarrow \det(A - 0 \cdot I) \neq 0 \Leftrightarrow (A - \lambda I) \neq 0 \Leftrightarrow \lambda = 0\)
is not a solution for \(p(\lambda) = 0\) \(\Leftrightarrow \lambda = 0\)
is \textbf{not} an eigenvalue of \(A\).

\begin{enumerate}
\def\labelenumi{\arabic{enumi}.}
\setcounter{enumi}{2}
\item
  \(A\) is invertible \(\Rightarrow\) \(\frac{1}{\lambda}\) is an
  eigenvalue of \(A^{- 1}\) corresponding to the eigenvector of
  \(\mathbf{v}\)
\end{enumerate}

\(A\mathbf{v} = \lambda\mathbf{v} \Rightarrow A^{- 1}\mathbf{v} = \frac{1}{\lambda}\mathbf{v}\)
(if \(A\) is invertible)

\emph{Proof.}

Firstly, \(\lambda \neq 0\), so \(\frac{1}{\lambda}\) makes sense.

\begin{align*}
A\mathbf{v} &= \lambda\mathbf{v} \\
A^{- 1}A\mathbf{v} &= A^{- 1}\lambda\mathbf{v} \\
\mathbf{v} &= A^{- 1}\mathbf{\lambda}\mathbf{v} \\
\mathbf{v} &= \lambda A^{- 1}\mathbf{v} \\
\frac{\mathbf{v}}{\lambda} &= A^{- 1}\mathbf{v} \\
\frac{1}{\lambda}\mathbf{v} &= A^{- 1}\mathbf{v}
\end{align*}




\end{proof}
\section{Diagonalization}

Let \(A,\ B \in M_{n \times n}\) (STRICTLY SQUARE MATRICES). \(A,\ B\)
are similar if:

\begin{enumerate}
\def\labelenumi{\arabic{enumi}.}
\item
  The determinants match (\(\det(A) = \det(B)\))
\item
  Both matrices are either invertible or not invertible (\(A\) is
  invertible iff \(B\) is invertible)
\item
  The rank and nullity of them match
  (\(\text{rank}(A) = \text{rank}(B)\);
  \(\text{nullity}(A) = \text{nullity}(B)\))
\item
  \(\det\left( A - \lambda I \right) = \det\left( B - \lambda I \right)\)
\item
  Same solutions for their characteristic polynomial meaning
  \(p\left( \lambda_{A} \right) = p\left( \lambda_{D} \right)\)
\item
  Same eigenvalues
\end{enumerate}

Proofs:

\begin{enumerate}
\def\labelenumi{\arabic{enumi}.}
\item
  --
\item
  Recall \(A\) is invertible if and only if \(\det(A) \neq 0\).
  \(\Rightarrow\) If \(A,\ B\) give you the same determinant, \(A\) is
  invertible if and only if \(B\) is invertible.
\item
  \(\text{rank}(A) + \text{nullity}(A) = n\), and
  \(\text{rank}(D) + \text{nullity}(D) = n\) (\(n\) is the column/row
  count of the matrix, and they have to at least be the same size) We
  need to show that the rank agrees with each other (which implies the
  nullities agree with each other).

  \begin{enumerate}
  \def\labelenumii{\alph{enumii}.}
  \item
    Show \(\text{rank}(A) = \text{rank}(D)\).
  \end{enumerate}
\end{enumerate}

\begin{quote}
\emph{Lemma:} \(\text{rank}\left( QM \right) = \text{rank}(M)\)
if \(Q\) is invertible (because we only care about the number of pivots
in the matrix -- reduce it to RREF and based off it we can count how
many pivots there are for the given matrix.)

You can review \(Q\) as a linear transformation, and if \(Q\) isn't
invertible it will squish everything into a lower dimension, removing
its bijectivity. That's why Q needs to be a bijective (isomorphic)
linear map \(\Rightarrow\) \textbf{Q PRESERVES ALL INFORMATION}. (Note
that for matrix multiplication, the order of application of linear maps
are from right to left)

If this is true, then
\(\text{rank}(B) = \text{rank}\left( P^{- 1}AP \right)\). By the
lemma,
\(\text{rank}\left( P^{- 1}AP \right) = \text{rank}\left( AP \right) = \text{rank}(A)\)
(as \(P\) is invertible).
\end{quote}

\begin{align*}
&= \det\left( P^{- 1}AP - P^{- 1} \cdot \lambda I \cdot P \right) = \det\left( P^{- 1}\left( A - \lambda I \right)P \right) \\
&= \frac{1}{\det(P)} \cdot \det\left( A - \lambda I \right) \cdot \det(P) = \det\left( A - \lambda I \right)
\end{align*}


\begin{enumerate}
\def\labelenumi{\arabic{enumi}.}
\setcounter{enumi}{3}
\item
  --
\item
  --
\end{enumerate}


\subsection{Diagonal matrices}

The determinant of a diagonal matrix is the product of its diagonal
entries. Computing the inverse matrix for the diagonal matrix is also
very easy (is the reciprocal of all entries in the diagonal matrix).

\begin{enumerate}
\def\labelenumi{\arabic{enumi}.}
\item
  Recall: A matrix \(A\) is diagonalizable if \(\exists\) an invertible
  \(n \times n\) matrix \(P\) such that \(D = P^{- 1}AP\) where
  \(D\) is a diagonal matrix.
\item
  If \(A\) and \(D\) are similar, then \(\exists P\) such that
  \(D = P^{- 1}AP\) or \(PDP^{- 1} = A\)

  \begin{enumerate}
  \def\labelenumii{\alph{enumii}.}
  \item
    \emph{Proof:}
    \(\det(D) = \det\left( P^{- 1}AP \right) = \det\left( P^{- 1} \right)\det(A)\det(P) = \frac{\det(P)}{\det(P)}\det(A) = \det(A)\)
  \end{enumerate}
\item
  We should know how to check if \(A\) is real-diagonalizable.
  \(p(\lambda) = 0\)
\end{enumerate}


\subsection{Check if a matrix is
diagonalizable}

Information about diagonalization: \(A \in M_{n \times n}\) (strictly
square)

\begin{enumerate}
\def\labelenumi{\arabic{enumi}.}
\item
  \(A\) is diagonalizable \(\Leftrightarrow\) \(A\) has \(n\) linearly
  independent eigenvectors.
\end{enumerate}

\(P(\lambda) = \det(A - \lambda I) = 0 \Rightarrow \lambda_{1}\cdots\lambda_{k}\)
eigenvalues for \(A\)

For \(E\left( \lambda_{1} \right)\), basis vector
\({\mathbf{v}}_{i}\) is a good representation for the eigenvector
corresponding to \(\lambda_{1}\),
\(\lambda_{1}{\mathbf{v}}_{i} = A{\mathbf{v}}_{i}\).
(\({\mathbf{v}}_{i}\) is a basic vector, which changes on which
eigenvalue you are checking) However, eigenspaces do not necessarily
contain one eigenvector, so you'll have to figure it out yourself.

The total number of eigenvector spaces must be equal to \(n\). Once
you've deduced \(n\) eigenvectors, you must verify that they are
linearly independent by putting them as column vectors in a matrix, then
row reducing it.

\begin{enumerate}
\def\labelenumi{\arabic{enumi}.}
\setcounter{enumi}{1}
\item
  If \(A \in M_{n \times n}\) and \(A\) has \(n\) \textbf{distinct}
  eigenvalues, \(\Rightarrow\) then \(A\) is diagonalizable. Write down
  the characteristic equation. \textbf{(This is a strictly an if then
  statement and not an iff)}
\end{enumerate}

\[p(\lambda) = 0,\deg\left( p(\lambda) \right) = \text{no.\ of\ distinct\ roots} \Rightarrow A\ \text{is\ diagonalizable}\]

\textbf{All distinct eigenvalues imply all the basis of all
eigenvectors} \({\mathbf{\mathbf{v}}}_{\mathbf{i}}\) \textbf{are
linearly independent.}


\subsection{The procedure of
diagonalization}

\begin{enumerate}
\def\labelenumi{\arabic{enumi}.}
\item
  Find
  \(\det(A - \lambda I) = 0 \Rightarrow \text{roots}\ \lambda_{1}\ldots\lambda_{n}\)
  (check if \(\text{A\ }\)is diagonalizable; all roots must be distinct)
\item
  Find
  \({\mathbf{v}}_{1},\ {\mathbf{v}}_{2}\ldots{\mathbf{v}}_{n}\)
  for \({\mathbf{v}}_{1}\ldots{\mathbf{v}}_{n}\) such that
  \(A{\mathbf{v}}_{i} = \lambda{\mathbf{v}}_{i}\forall i\)
\item
  Use them to construct matrix \(P = \begin{bmatrix}
   \mid & \mid & \mid & \mid & \mid \\
  {\mathbf{v}}_{1} & {\mathbf{v}}_{2} & {\mathbf{v}}_{3} & \cdots & {\mathbf{v}}_{n} \\
   \mid & \mid & \mid & \mid & \mid \\
  \end{bmatrix}\) and \(D = \begin{bmatrix}
  \lambda_{1} & 0 & 0 & 0 \\
  0 & \lambda_{2} & 0 & 0 \\
  0 & 0 & \ddots & 0 \\
  0 & 0 & 0 & \lambda_{n} \\
  \end{bmatrix}\)
\item
  \(P^{- 1} \cdot A \cdot P = D\)
\end{enumerate}

Why is this true? \emph{Proof.}
\(P^{- 1}AP = D \Leftrightarrow AP = PD\) (pre-multiply \(P\) on both
sides)

Verify if \(AP = PD\). Compute LHS and RHS.

\[{LHS = AP }{= A\begin{bmatrix}
 \mid & \mid & \mid & \mid & \mid \\
{\mathbf{v}}_{1} & {\mathbf{v}}_{2} & {\mathbf{v}}_{3} & \cdots & {\mathbf{v}}_{n} \\
 \mid & \mid & \mid & \mid & \mid \\
\end{bmatrix} }{= \left\lbrack A\begin{matrix}
 \mid & \mid & \mid & \mid & \mid \\
{\mathbf{v}}_{1} & A{\mathbf{v}}_{2} & A{\mathbf{v}}_{3} & \cdots & A{\mathbf{v}}_{n} \\
 \mid & \mid & \mid & \mid & \mid \\
\end{matrix} \right\rbrack} \]

\[A{\mathbf{v}}_{i} = \lambda_{i}{\mathbf{v}}_{i}\]

\[{}{= \left\lbrack \lambda_{1}\begin{matrix}
 \mid & \mid & \mid & \mid & \mid \\
{\mathbf{v}}_{1} & \lambda_{2}{\mathbf{v}}_{2} & \lambda_{3}{\mathbf{v}}_{3} & \cdots & \lambda_{n}{\mathbf{v}}_{n} \\
 \mid & \mid & \mid & \mid & \mid \\
\end{matrix} \right\rbrack }{= \begin{bmatrix}
 \mid & \mid & \mid & \mid & \mid \\
{\mathbf{v}}_{1} & {\mathbf{v}}_{2} & {\mathbf{v}}_{3} & \cdots & {\mathbf{v}}_{n} \\
 \mid & \mid & \mid & \mid & \mid \\
\end{bmatrix} \cdot \begin{bmatrix}
\lambda_{1} & 0 & 0 & 0 \\
0 & \lambda_{2} & 0 & 0 \\
0 & 0 & \ddots & 0 \\
0 & 0 & 0 & \lambda_{n} \\
\end{bmatrix} }{= PD = RHS} \]

Example: \(A = \begin{bmatrix}
2 & 0 & 0 \\
1 & 3 & 0 \\
 - 3 & 5 & 4 \\
\end{bmatrix}\). Check if \(A\) is diagonalizable. If yes, find
\(D,\ P\) such that \(P^{- 1}AP = D\).

\begin{enumerate}
\def\labelenumi{\arabic{enumi}.}
\item
  Find all eigenvalues.
\end{enumerate}

\[\det\left( \begin{bmatrix}
2 - \lambda & 0 & 0 \\
1 & 3 - \lambda & 0 \\
 - 3 & 5 & 4 - \lambda \\
\end{bmatrix} \right) = (\lambda - 2)(\lambda - 3)(\lambda - 4) = 0\]

The eigenvalues: \(\lambda_{1} = 2,\ \lambda_{2} = 3,\ \lambda_{3} = 4\)

For \(\lambda = 2\), \(A - 2I = \begin{bmatrix}
0 & 0 & 0 \\
1 & 1 & 0 \\
3 & 5 & 2 \\
\end{bmatrix};\text{null}(A - \lambda I) = sp\left( \begin{bmatrix}
1 \\
 - 1 \\
4 \\
\end{bmatrix} \right)\). Your basis is\(\begin{bmatrix}
1 \\
 - 1 \\
4 \\
\end{bmatrix} = {\mathbf{v}}_{1}\);
\(\lbrack A\rbrack\begin{bmatrix}
1 \\
 - 1 \\
4 \\
\end{bmatrix} = \lambda\begin{bmatrix}
1 \\
 - 1 \\
4 \\
\end{bmatrix}\) where \(\lambda = 2\).

For \(\lambda = 3\), \(A - 3I = \begin{bmatrix}
 - 1 & 0 & 0 \\
1 & 0 & 0 \\
 - 3 & 5 & 1 \\
\end{bmatrix},\text{null}(A - 3I) = sp\left( \begin{bmatrix}
0 \\
1 \\
 - 5 \\
\end{bmatrix} \right)\). Your basis is \(\begin{bmatrix}
0 \\
1 \\
 - 5 \\
\end{bmatrix} = {\mathbf{v}}_{2}\).

For \(\lambda = 4,\ A - 4I = \begin{bmatrix}
 - 2 & 0 & 0 \\
 - 1 & - 1 & 0 \\
 - 3 & 5 & 0 \\
\end{bmatrix},\text{null}(A - 4I) = sp\left( \begin{bmatrix}
0 \\
0 \\
1 \\
\end{bmatrix} \right)\). Your basis is \(\begin{bmatrix}
0 \\
0 \\
1 \\
\end{bmatrix} = {\mathbf{v}}_{3}\).

\({\mathbf{v}}_{1},\ {\mathbf{v}}_{2},\ {\mathbf{v}}_{3}\)
should be linearly independent, as the eigenvalues are distinct.

\[P = \begin{bmatrix}
 \mid & \mid & \mid \\
{\mathbf{v}}_{1} & {\mathbf{v}}_{2} & {\mathbf{v}}_{3} \\
 \mid & \mid & \mid \\
\end{bmatrix} = \begin{bmatrix}
1 & 0 & 0 \\
 - 1 & 1 & 0 \\
4 & - 5 & 1 \\
\end{bmatrix}\]

Fact: Let
\({\mathbf{v}}_{1},\ {\mathbf{v}}_{2},\ \ldots,\ {\mathbf{v}}_{k}\)
be eigenvectors corresponding to distinct eigenvalues
\(\lambda_{1}\ldots\lambda_{k}\). Then,
\({\mathbf{v}}_{1}\ldots{\mathbf{v}}_{k}\) are linear
independent \(\Rightarrow P = \begin{bmatrix}
 \mid & \mid & \mid & \mid & \mid \\
{\mathbf{v}}_{1} & {\mathbf{v}}_{2} & {\mathbf{v}}_{3} & \cdots & {\mathbf{v}}_{k} \\
 \mid & \mid & \mid & \mid & \mid \\
\end{bmatrix} \Rightarrow\) \(P\) is invertible. Then

\[P^{- 1}AP = D = \begin{bmatrix}
2 & 0 & 0 \\
0 & 3 & 0 \\
0 & 0 & 4 \\
\end{bmatrix}\]

Where \(A\) is similar to \(D\).

\[\det(A) = \det(D) = 2 \cdot 3 \cdot 4 = 24\]

\begin{definition}
Let \(A\) be an \(n \times n\) matrix with
characteristic polynomial

\[p(\lambda) = \left( \lambda - \lambda_{1} \right)^{\alpha_{1}}\left( \lambda - \lambda_{2} \right)^{\alpha_{2}}\cdots\left( \lambda - \lambda_{i} \right)^{\alpha i}\]
\end{definition}

\(p(\lambda) = 0 \Rightarrow \lambda_{1},\ \lambda_{2},\ \cdots,\ \lambda_{i}\ \)with
order \(\alpha_{1}\cdots\alpha_{i}\).

Algebraic multiplicity of \(\lambda_{i}\) is \(\alpha_{i}\)

Geometric multiplicity of \(\lambda_{i}\) is
\(\dim\left( E_{\lambda_{i}} \right) = \beta_{i}\)

\textbf{Fact:} For every \(\lambda_{i}:\) If
\(\text{alg}\left( \lambda_{i} \right) = geo\left( \lambda_{i} \right) \Rightarrow A\)
is diagonalizable \textbf{(Matching algebraic and geometric multiplicity
FOR ALL eigenvalues will guarantee you are diagonalizable matrix)}


\subsection{The diagonalization
process}

Let \(A \in M_{n \times n}\).

\begin{enumerate}
\def\labelenumi{\arabic{enumi}.}
\item
  Find the characteristic polynomial
  \(p(\lambda) = \det(A - \lambda I)\)
\item
  Factor.
  \(p(\lambda) = \left( \lambda - \lambda_{1} \right)^{n_{1}}\left( \lambda - \lambda_{2} \right)^{n_{2}}\cdots\left( \lambda - \lambda_{i} \right)^{n_{i}}\)
\end{enumerate}

\begin{align*}
\text{alg}\left( \lambda_{i} \right) &= n_{i} \\
\text{geo}\left( \lambda_{i} \right) &= \dim\left( E_{\lambda_{i}} \right) = \dim\left( \text{null}\left( A - \lambda_{i}I \right) \right) = \text{nullity}\left( A - \lambda_{i}I \right)
\end{align*}


To figure \(\text{null}\left( A - \lambda_{i}I \right)\),
\(\left( A - \lambda_{i}I \right)\mathbf{x} = \mathbf{0}\) and
figure out how many free variables, which tells us the dimension of the
nullspace. This tells us the dimension of the eigenspace.

Once you know the dimension for the eigenspace, you will know the
geometric multiplicity.

\begin{enumerate}
\def\labelenumi{\arabic{enumi}.}
\setcounter{enumi}{2}
\item
  Check if
  \(\text{alg}\left( \lambda_{i} \right) = \text{geo}\left( \lambda_{i} \right)\)
  for every \(\lambda_{i}\). If that is the case the matrix is
  diagonalizable.
\item
  If \(A\) is diagonalizable, then you can diagonalize \(A\) by the
  following:
\end{enumerate}

\[A\sim D = \begin{bmatrix}
\lambda_{1} & 0 & 0 & 0 \\
0 & \lambda_{2} & 0 & 0 \\
0 & 0 & \ddots & 0 \\
0 & 0 & 0 & \lambda_{n} \\
\end{bmatrix}\]

An example:

If \(A\) is an \(n \times n\) matrix and \(A\) is diagonalizable. Then
there is a shortcut to compute \(A^{k}\).

Solution: \(P^{- 1}AP = D = \begin{bmatrix}
\lambda_{1} & 0 & 0 \\
0 & \ddots & 0 \\
0 & 0 & \lambda_{n} \\
\end{bmatrix}\)

\begin{align*}
A &= PDP^{- 1} \\
A^{k} &= \left( PDP^{- 1} \right)\left( PDP^{- 1} \right)\cdots\left( PDP^{- 1} \right)\ \text{k\ times} = P \cdot D^{k} \cdot P^{- 1}
\end{align*}


Note that because matrix multiplication is associative you have a bunch
of \(P^{- 1}P\) cases.

Then \(D^{k} = \begin{bmatrix}
\lambda_{1}^{k} & 0 & 0 \\
0 & \ddots & 0 \\
0 & 0 & \lambda_{n}^{k} \\
\end{bmatrix}\)

\emph{When you mention an eigenvector, it must correspond to an
eigenvalue}

\emph{The rank-nullity theorem will appear in the final exam.}

Determine whether matrix \(A = \begin{bmatrix}
0 & 0 & - 2 \\
1 & 2 & 1 \\
1 & 0 & 3 \\
\end{bmatrix}\) is real diagonalizable.

\[P^{- 1}AP = D = \begin{bmatrix}
\lambda_{1} & 0 & 0 \\
0 & \lambda_{2} & 0 \\
0 & 0 & \lambda_{3} \\
\end{bmatrix}\]

\begin{align*}
p(A) &= \det(A - \lambda I) \\
(a - \lambda)(b - \lambda) - cd &= 0 \\
\Delta &< 0 \Rightarrow \text{Can't\ be\ real\ diagonalized}
\end{align*}


Solution:

\[p(\lambda) = \det(A - \lambda I) = \left| \begin{matrix}
 - \lambda & 0 & - 2 \\
1 & - 2 - \lambda & 1 \\
1 & 0 & 3 - \lambda \\
\end{matrix} \right| = (2 - \lambda)\left( \lambda^{2} - 3\lambda + 2 \right) = (2 - \lambda)(\lambda - 2)(\lambda - 1)\]

\[\lambda_{1} = 2,\ \lambda_{2} = 2,\ \lambda_{3} = 1\]

You have a repeated eigenvalue. This means the eigenvalue 2 repeats
twice, meaning the algebraic multiplicity for \(\lambda = 2\) to be two.
This means we expect the geometric multiplicity to be two as well.

Meaning \(\dim\left( E_{\lambda = 2} \right) = 2\). If this were not the
case, \(\text{alg}(A) \neq \text{geo}(A) \Rightarrow A\) wouldn't be
diagonalizable.

Checking eigenspaces: For \(\lambda = 2,\text{null}(A - \lambda I)\)
(the nullspace is the solution space of
\((A - \lambda I)\mathbf{x} = \mathbf{0}\))

\[A - \lambda I = \begin{bmatrix}
 - 2 & 0 & - 2 \\
1 & 0 & 1 \\
1 & 0 & 1 \\
\end{bmatrix}\sim\begin{bmatrix}
1 & 0 & 1 \\
0 & 0 & 0 \\
0 & 0 & 0 \\
\end{bmatrix} \Rightarrow \text{nullity}(A - \lambda I) = 2\]

Based on this I can tell there are two free variables, meaning
\(\dim\left( \mathbf{x} \right) = 2 \Rightarrow \dim\left( E_{\lambda = 2} \right) = 2\).

For \(\lambda = 1\)

\[A - \lambda I = \begin{bmatrix}
 - 1 & 0 & - 2 \\
1 & 1 & 1 \\
1 & 0 & 2 \\
\end{bmatrix}\sim\begin{bmatrix}
1 & 0 & 2 \\
0 & 1 & - 1 \\
0 & 0 & 0 \\
\end{bmatrix}\]

Two columns have a pivot \(\Rightarrow 1\) free variable
\(\Rightarrow \dim\left( \mathbf{x} \right) = 1 \Rightarrow \dim\left( E_{\lambda = 1} \right) = 1 \Rightarrow \text{geo}\left( E_{\lambda = 1} \right) = 1\)

The eigenvalue only appears one time during the factorization of the
characteristic polynomial, so the algebraic multiplicity of this
eigenvalue is 1. Therefore,
\(\text{alg}(\lambda = 1) = \text{geo}(\lambda = 1)\).

\(\Rightarrow A\) is diagonalizable; \(\Rightarrow A\) is similar to
\(D = \begin{bmatrix}
2 & 0 & 0 \\
0 & 2 & 0 \\
0 & 0 & 1 \\
\end{bmatrix}\), \(\det(D) = 2 \cdot 2 \cdot 1 = 4\)

\emph{What can we get out of this?}

\begin{enumerate}
\def\labelenumi{\arabic{enumi}.}
\item
  \(\det(A) = \det(D) = 2 \cdot 2 \cdot 1 = 4\) \emph{As A is similar to
  D, their determinants match.}
\item
  Is \(A\) invertible? Yes. \(A\) is invertible as \(D\) is invertible
  \emph{(and also because we just computed the determinant of} \(A\)
  \emph{already and it is nonzero)}
\item
  \(T\left( \mathbf{x} \right) = A\mathbf{x}:\mathbb{R}^{3} \rightarrow \mathbb{R}^{3}\)
  (And this is bijective because \(A\) is invertible -
  \(A\sim I;A\mathbf{x} = \mathbf{0} \Rightarrow \mathbf{x} = \mathbf{0}\).)
\item
  \(\text{nullity}(A) = 0\); \(A\) is full rank as
  \(\text{rank}(A) = 3\).
\end{enumerate}

Rank-nullity theorem: \(\text{nullity}(A) + \text{rank}(A) = n\);
\(0 + 3 = 3\)

\begin{enumerate}
\def\labelenumi{\arabic{enumi}.}
\setcounter{enumi}{4}
\item
  \(A\mathbf{x} = \mathbf{b}\) will always have a unique
  solution.
\end{enumerate}


\subsection{\texorpdfstring{Every symmetric matrix is diagonalizable
\(\mathbf{A}^{\mathbf{T}}\mathbf{= A}\)}{Every symmetric matrix is diagonalizable \textbackslash mathbf\{A\}\^{}\{\textbackslash mathbf\{T\}\}\textbackslash mathbf\{= A\}}}

Example: \(A = \begin{bmatrix}
3 & - 3 & - 3 \\
 - 3 & 3 & - 3 \\
 - 3 & - 3 & 3 \\
\end{bmatrix}\); \(A^{T} = A \Rightarrow A\) is symmetric
\(\Rightarrow A\) is diagonalizable \(\Rightarrow A\) is similar to
\(D = \begin{bmatrix}
\lambda_{1} & 0 & 0 \\
0 & \lambda_{2} & 0 \\
0 & 0 & \lambda_{3} \\
\end{bmatrix}\)

(Excuse: from lecture, every symmetric matrix is diagonalizable.)

Computing the eigenvalues for \(A\):

\(p(\lambda) = \det(A - \lambda I) = \left| \begin{matrix}
3 - \lambda & - 3 & - 3 \\
 - 3 & 3 - \lambda & - 3 \\
 - 3 & - 3 & 3 - \lambda \\
\end{matrix} \right| = (6 - \lambda)(\lambda + 3)(\lambda - 6) = 0 \Rightarrow\)
Eigenvalues: \(\lambda_{1} = 6,\ \lambda_{2} = - 3,\ \lambda_{3} = 6\)

\emph{In this case, you have one eigenvalue repeating twice, and one
eigenvalue appearing only once. Normally, you would have to verify if}
\(\text{alg}(\lambda = 6) = \text{geo}(\lambda = 6) = 2\) \emph{and the
one for 3, but because the matrix is diagonalizable, you don't need to.}

\begin{align*}
&\dim\left( E_{\lambda = 6} \right) = 2 \\
A - 6I &= \begin{bmatrix} - 3 & - 3 & - 3 \\ - 3 & - 3 & - 3 \\ - 3 & - 3 & - 3 \\\end{bmatrix}\sim\begin{bmatrix}1 & 1 & 1 \\0 & 0 & 0 \\0 & 0 & 0 \\\end{bmatrix};\text{nullity}(A - 6I) = 2
\end{align*}


\[x_{3} = s;x_{2} = t;\mathbf{x} = \begin{bmatrix}
x_{1} \\
x_{2} \\
x_{3} \\
\end{bmatrix} = \begin{bmatrix}
 - t - s \\
t \\
s \\
\end{bmatrix} = t\begin{bmatrix}
 - 1 \\
1 \\
0 \\
\end{bmatrix} + s\begin{bmatrix}
 - 1 \\
0 \\
1 \\
\end{bmatrix}\]

\begin{align*}
&\dim\left( E_{\lambda = 1} \right) = 1 \\
&E_{\lambda = 1} = sp\left( \begin{bmatrix}1 \\1 \\1 \\\end{bmatrix} \right)
\end{align*}


Note: To make \(P\) nicely, exactly 3 linearly independent eigenvectors
must appear in your solutions. In this case, we have exactly 3, so we
can use them to make up \(P\):

\[P = \begin{bmatrix}
1 & 1 & 0 \\
1 & 0 & 1 \\
1 & - 1 & - 1 \\
\end{bmatrix}\]

Composed from the column vectors of
\({\mathbf{v}}_{1},\ {\mathbf{v}}_{2},\ {\mathbf{v}}_{3}\).
The order of what vectors from the eigenvalues you put does not matter,
but \textbf{it will affect} the column ordering of \(D\). For example,
if \({\mathbf{v}}_{1}\) is an eigenvector of \(\lambda = 1\), \(1\)
has to be put on the top left of the matrix.

\[P^{- 1}AP = D = \begin{bmatrix}
 - 3 & 0 & 0 \\
0 & 6 & 0 \\
0 & 0 & 6 \\
\end{bmatrix}\]


\subsection{\texorpdfstring{ Justifying
order}{ Justifying order}}

\[{A \cdot \begin{bmatrix}
 \mid & \mid & \mid \\
{\mathbf{v}}_{3} & {\mathbf{v}}_{2} & {\mathbf{v}}_{1} \\
 \mid & \mid & \mid \\
\end{bmatrix} }{= \begin{bmatrix}
 \mid & \mid & \mid \\
{\mathbf{v}}_{3} & {\mathbf{v}}_{2} & {\mathbf{v}}_{1} \\
 \mid & \mid & \mid \\
\end{bmatrix}\begin{bmatrix}
d_{1} & 0 & 0 \\
0 & d_{2} & 0 \\
0 & 0 & d_{3} \\
\end{bmatrix} }{= \begin{bmatrix}
 \mid & \mid & \mid \\
{A\mathbf{v}}_{3} & A{\mathbf{v}}_{2} & A{\mathbf{v}}_{1} \\
 \mid & \mid & \mid \\
\end{bmatrix} }{= \left\lbrack d_{1}\begin{matrix}
 \mid & \mid & \mid \\
{\mathbf{v}}_{3} & d_{2}{\mathbf{v}}_{2} & {d_{3}\mathbf{v}}_{1} \\
 \mid & \mid & \mid \\
\end{matrix} \right\rbrack} \]

\emph{Order of} \({\mathbf{v}}_{i}\) \emph{will determine the order
of eigenvalues you put in D.}


\subsection{Property of matrix exponentiation with
diagonalization:}

\emph{Theorem.}

\begin{align*}
A &= PDP^{- 1} \\
&\Rightarrow A^{k}  \\
&= \left( PDP^{- 1} \right)\left( PDP^{- 1} \right)\left( PDP^{- 1} \right)\cdots\left( PDP^{- 1} \right)\text{\ k\ }\text{times} \\
\Rightarrow A^{k} &= PD\left( P^{- 1}P \right)D\left( P^{- 1}P \right)D\left( P^{- 1}P \right)\cdots\left( P^{- 1}P \right)DP^{- 1} \\
\Rightarrow A^{k} &= PD^{k}P^{- 1}
\end{align*}

\end{document}
